\documentclass[11pt,a4paper]{article}
\usepackage[utf8]{inputenc}
% lmodern must precede fontenc: without a scalable Type1 font the T1 encoding
% falls back to bitmap fonts, and microtype's expansion then aborts the run
% with "auto expansion is only possible with scalable fonts".
\usepackage{lmodern}
\usepackage[T1]{fontenc}
\usepackage[margin=1in]{geometry}
\usepackage{amsmath,amssymb}
\usepackage{graphicx}
\usepackage{microtype}
\usepackage[hidelinks,breaklinks]{hyperref}
\usepackage{url}
\urlstyle{same}

% Generated by code/tools/md_to_tex.py from the SurveyForge Markdown output.
% Citation numbers match the source .md exactly.

\title{Comprehensive Survey on Visual Adversarial Attacks and Defenses in the Physical World}
\author{Generated by SurveyForge}
\date{}

\begin{document}
\maketitle
\tableofcontents
\newpage



\section{Introduction}


The importance of visual adversarial attacks and defenses in the physical world cannot be overstated, as they have significant implications for security, safety, and reliability in various applications such as autonomous vehicles, surveillance systems, and robotics \cite{doi:10.24963/ijcai.2021/694,doi:10.1007/s11633-019-1211-x}. These attacks involve manipulating inputs to machine learning models in ways that are imperceptible to humans but can cause the models to make incorrect predictions \cite{doi:10.1145/3128572.3140444}. In the physical world, such attacks can be realized through various means, including printed adversarial patches or camouflage, which can be used to deceive object detectors and classifiers \cite{arxiv:1707.08945,arxiv:1712.09665}. The potential consequences of these attacks are far-reaching, ranging from compromising the safety of autonomous vehicles to undermining the reliability of surveillance systems \cite{doi:10.1007/978-3-030-10925-7_4,arxiv:1906.11897}.

Recent studies have demonstrated the feasibility of physical adversarial attacks in various scenarios, including attacking stop signs \cite{arxiv:1707.08945} and object detectors \cite{doi:10.1007/978-3-030-58574-7_41}. These attacks have been shown to be effective even in the presence of various environmental factors, such as lighting and viewpoint changes \cite{doi:10.1109/tnnls.2023.3321432,doi:10.1109/tpami.2024.3430860}. The development of robust defenses against these attacks is therefore crucial, and several approaches have been proposed, including adversarial training, defensive distillation, and certification methods \cite{doi:10.1109/iccv48922.2021.01542,arxiv:2003.06693,arxiv:1902.06705}.

Despite the progress made in understanding and mitigating visual adversarial attacks, several challenges remain, including the need for more robust and generalizable defenses, as well as the development of evaluation metrics and methodologies that can accurately assess the effectiveness of these attacks and defenses \cite{arxiv:1902.06705,doi:10.1109/access.2022.3216291}. Furthermore, the emergence of new technologies, such as 3D printing and augmented reality, is likely to introduce new vulnerabilities and challenges in the physical world \cite{doi:10.1109/tpami.2024.3430860,doi:10.1016/j.ins.2023.03.084}. As the field continues to evolve, it is essential to stay ahead of these challenges and develop innovative solutions that can ensure the security, safety, and reliability of machine learning systems in the physical world \cite{doi:10.1016/j.eng.2019.12.012,arxiv:1810.00069}.

The development of effective defenses against visual adversarial attacks requires a deep understanding of the underlying mechanisms and vulnerabilities of machine learning models \cite{doi:10.24963/ijcai.2021/694,doi:10.1007/s11633-019-1211-x}. Several approaches have been proposed, including adversarial training, which involves training models on adversarial examples to improve their robustness \cite{doi:10.1109/iccv48922.2021.01542,arxiv:2003.06693}. Other approaches, such as defensive distillation and certification methods, have also been shown to be effective in improving the robustness of machine learning models \cite{arxiv:2002.07405,arxiv:2003.06693}. However, the development of more robust and generalizable defenses remains an open challenge, and further research is needed to address this issue \cite{arxiv:1902.06705,doi:10.1109/access.2022.3216291}.

In conclusion, visual adversarial attacks and defenses in the physical world are a critical area of research, with significant implications for security, safety, and reliability in various applications. While significant progress has been made in understanding and mitigating these attacks, several challenges remain, including the need for more robust and generalizable defenses, as well as the development of evaluation metrics and methodologies that can accurately assess the effectiveness of these attacks and defenses \cite{doi:10.24963/ijcai.2021/694,doi:10.1007/s11633-019-1211-x}. As the field continues to evolve, it is essential to stay ahead of these challenges and develop innovative solutions that can ensure the security, safety, and reliability of machine learning systems in the physical world \cite{doi:10.1016/j.eng.2019.12.012,arxiv:1810.00069}. Future research directions include the development of more robust and generalizable defenses, as well as the exploration of new application domains and the investigation of new attack and defense strategies \cite{doi:10.1016/j.eng.2019.12.012,arxiv:1810.00069}.


\section{Fundamentals of Visual Adversarial Attacks}



\subsection{Introduction to Visual Adversarial Attacks}


The concept of visual adversarial attacks has gained significant attention in recent years due to its potential to compromise the security and reliability of machine learning models in the physical world. As defined in \cite{doi:10.1145/3128572.3140444}, visual adversarial attacks refer to the process of crafting inputs that are specifically designed to mislead machine learning models, while being imperceptible to human observers. These attacks can be categorized into different types, including patch-based, object-based, and scene-based attacks, each with its unique characteristics and challenges \cite{arxiv:1707.08945}.

Patch-based attacks, for instance, involve adding a small perturbation to the input image, which can be in the form of a sticker or a patch, to fool the machine learning model \cite{arxiv:1712.09665}. These attacks have been shown to be effective in various scenarios, including object detection and image classification \cite{arxiv:1906.11897}. Object-based attacks, on the other hand, involve manipulating the object itself, rather than adding a patch, to create an adversarial example \cite{arxiv:1807.07769}. Scene-based attacks, as the name suggests, involve manipulating the entire scene, rather than a specific object or patch, to create an adversarial example \cite{doi:10.24963/ijcai.2021/694}.

The importance of understanding and mitigating visual adversarial attacks cannot be overstated. As machine learning models become increasingly ubiquitous in various applications, including security-critical systems, the potential consequences of these attacks can be severe \cite{doi:10.1109/access.2018.2807385}. For instance, an adversarial attack on a self-driving car's object detection system can have catastrophic consequences \cite{doi:10.1007/978-3-030-10925-7_4}. Therefore, it is essential to develop effective defense mechanisms to detect and mitigate these attacks \cite{arxiv:1810.00069}.

Recent studies have proposed various defense mechanisms, including detection-based defenses, mitigation strategies, and robustness enhancement techniques \cite{doi:10.1049/cit2.12028}. Detection-based defenses involve detecting adversarial examples before they can cause harm, while mitigation strategies involve reducing the impact of these attacks. Robustness enhancement techniques, on the other hand, involve improving the robustness of machine learning models to adversarial attacks \cite{arxiv:2003.06693}. These defense mechanisms have shown promising results in various scenarios, including image classification and object detection \cite{doi:10.1016/j.cose.2020.102092}.

Despite the progress made in developing defense mechanisms, visual adversarial attacks remain a significant challenge. The complexity of these attacks, combined with the limited understanding of the underlying mechanisms, makes it difficult to develop effective defenses \cite{doi:10.1007/s11633-019-1211-x}. Furthermore, the emergence of new attack methods, such as 3D printing-based attacks and drone-based attacks, has further complicated the situation. Therefore, it is essential to continue researching and developing new defense mechanisms to stay ahead of these evolving threats.

In conclusion, visual adversarial attacks pose a significant threat to the security and reliability of machine learning models in the physical world. Understanding and mitigating these attacks is crucial to ensuring the safe and reliable operation of these models. While significant progress has been made in developing defense mechanisms, further research is needed to stay ahead of the evolving threats and to develop more effective defenses \cite{doi:10.1109/access.2022.3216291}. As noted in \cite{doi:10.58496/adsa/2024/011}, the development of more robust and generalizable defense mechanisms is essential to ensuring the security and reliability of machine learning models in the physical world. By continuing to research and develop new defense mechanisms, we can ensure the safe and reliable operation of machine learning models and prevent the potential consequences of visual adversarial attacks \cite{doi:10.1109/lra.2022.3189783}.


\subsection{Methods for Generating Visual Adversarial Examples}


The generation of visual adversarial examples is a crucial aspect of understanding and mitigating the vulnerabilities of deep learning models in the physical world, as discussed in the previous section on the importance of developing effective defense mechanisms against visual adversarial attacks. These examples are crafted to deceive machine learning models into misclassifying or misinterpreting the input, and they can be created using various methods, including optimization-based and generative model-based approaches. At the forefront of these methods are techniques such as the Fast Gradient Sign Method (FGSM) \cite{doi:10.24963/ijcai.2021/694} and the Projected Gradient Descent (PGD) \cite{arxiv:1902.06705}, which involve iteratively updating the input to maximize the loss function of the target model. These methods are effective in generating adversarial examples but can result in perturbations that are noticeable to the human eye, as demonstrated in \cite{arxiv:1712.09665}, which shows how optimization-based methods can be used to create adversarial patches that can fool object detectors.

On the other hand, generative model-based approaches, such as Generative Adversarial Networks (GANs) \cite{doi:10.1007/s11633-019-1211-x} and Variational Autoencoders (VAEs) \cite{doi:10.1109/tnnls.2025.3561225}, leverage the power of generative models to produce more realistic and subtle adversarial examples. These models can learn the distribution of the input data and generate perturbations that are more natural and less perceptible, as seen in \cite{doi:10.1007/978-3-030-58574-7_41}, which uses a generative model to create textured patches that can fool object detectors without being noticeable. The choice between optimization-based and generative model-based approaches depends on the specific use case and the desired level of stealth and effectiveness, highlighting the need for a comprehensive understanding of the strengths and limitations of each approach.

Recent studies have also explored the use of hybrid approaches that combine the strengths of optimization-based and generative model-based methods, such as \cite{arxiv:1806.02299}, which proposes a hybrid approach that uses a generative model to create the initial patch and then refines it using an optimization-based method. This approach can result in more effective and stealthy adversarial patches, underscoring the potential for hybrid methods to enhance the security and reliability of deep learning models. The development of methods for generating visual adversarial examples has significant implications for the security and reliability of deep learning models in the physical world, particularly as these models become increasingly ubiquitous in safety-critical applications, such as autonomous vehicles and surveillance systems.

As research in this area continues to evolve, it is essential to consider the properties of adversarial examples, including imperceptibility, transferability, and robustness, which will be discussed in the following section. By understanding these properties and developing effective methods for generating visual adversarial examples, researchers can identify potential weaknesses in deep learning models and develop more robust and secure architectures. Ultimately, the study of visual adversarial examples has the potential to significantly improve the security and reliability of deep learning models in the physical world, and it is an area that warrants further research and attention \cite{doi:10.1109/lra.2022.3189783}. Future directions in this field may include the exploration of new generative models and hybrid approaches, as well as the development of more effective defense mechanisms against visual adversarial attacks \cite{doi:10.1016/j.eng.2019.12.012}.


\subsection{Properties of Adversarial Examples}


The properties of adversarial examples are crucial in understanding their implications for attack effectiveness and defense strategies. At the forefront of these properties are imperceptibility, transferability, and robustness. Imperceptibility refers to the degree to which an adversarial example is indistinguishable from a legitimate input \cite{doi:10.1145/3128572.3140444}. This property is often measured using metrics such as peak signal-to-noise ratio (PSNR) and structural similarity index (SSIM) to evaluate the similarity between adversarial and original examples \cite{arxiv:1707.07397}. Transferability, on the other hand, pertains to the ability of adversarial examples to fool multiple models, including those with different architectures and training data \cite{doi:10.1109/cvpr.2019.00284}. This property has significant implications for black-box attacks, where the attacker has no knowledge of the target model's architecture or parameters.

Robustness is another critical property, referring to the ability of adversarial examples to withstand transformations and perturbations \cite{doi:10.1109/tip.2021.3092822}. Adversarial examples that are robust can maintain their effectiveness even when subjected to various types of noise or transformations, making them more potent in real-world scenarios. The robustness of adversarial examples can be enhanced through techniques such as adversarial training and the use of generative models \cite{arxiv:1805.06605}. Furthermore, the concept of natural adversarial examples \cite{doi:10.1007/s13042-021-01435-0} highlights that adversarial examples can occur naturally, without the need for deliberate crafting, which underscores the complexity of defending against such examples.

The interplay between these properties and the effectiveness of defense strategies is complex. For instance, defenses that focus on detecting adversarial examples based on their imperceptibility may be less effective against robust adversarial examples \cite{doi:10.1609/aaai.v32i1.11828}. Similarly, defenses that rely on transferability may be vulnerable to targeted attacks designed to evade such defenses \cite{doi:10.24963/ijcai.2019/134}. Therefore, understanding the properties of adversarial examples is essential for developing effective defense mechanisms.

Recent studies have also explored the use of evolutionary algorithms \cite{doi:10.1186/s41074-019-0053-3} and Bayesian optimization \cite{doi:10.1609/aaai.v34i07.6918} to generate adversarial examples with improved properties. These approaches have shown promise in creating adversarial examples that are not only effective but also robust against various defense mechanisms. Moreover, the application of adversarial examples in domains beyond image classification, such as object detection \cite{doi:10.24963/ijcai.2019/134} and video recognition \cite{doi:10.1609/aaai.v34i07.6918}, further emphasizes the need for comprehensive understanding and mitigation strategies.

In conclusion, the properties of adversarial examples, including imperceptibility, transferability, and robustness, play a crucial role in determining their effectiveness and the challenges they pose to defense strategies. As research in this area continues to evolve, it is essential to consider these properties in the development of more robust and generalizable defenses. Future directions may include exploring new properties of adversarial examples, such as their interpretability \cite{doi:10.1007/s00521-021-06812-y}, and leveraging techniques from other domains to enhance the security of machine learning models against adversarial attacks. By synthesizing insights from various studies, researchers can pave the way for more secure and reliable machine learning systems.


\subsection{Impact of Adversarial Examples on Deep Learning Models}


The impact of adversarial examples on deep learning models is a critical concern, as these models are increasingly being deployed in real-world applications, including image classification, object detection, and image segmentation. As discussed earlier, the properties of adversarial examples, such as imperceptibility, transferability, and robustness, play a crucial role in determining their effectiveness and the challenges they pose to defense strategies. Adversarial examples are inputs to a model that are specifically designed to cause the model to make a mistake, and they can be crafted to target specific models or to be transferable across different models \cite{doi:10.1145/3128572.3140444}. The existence of adversarial examples raises questions about the robustness and reliability of deep learning models, and it has significant implications for their use in safety-critical and security-sensitive applications \cite{arxiv:1511.07528}.

In the context of image classification, adversarial examples can cause a model to misclassify an image, even if the image is only slightly modified \cite{doi:10.1201/9781351251389-8}. This can have serious consequences, such as in self-driving cars or medical diagnosis, where incorrect classifications can have life-or-death consequences. Moreover, adversarial examples can be designed to be transferable across different models, making them a potent attack vector \cite{doi:10.1109/cvpr.2019.00284}. Researchers have proposed various methods to generate adversarial examples, including optimization-based methods and generative model-based methods \cite{arxiv:1707.07397}, and these methods have been shown to be effective in generating adversarial examples that can fool state-of-the-art models.

The impact of adversarial examples is not limited to image classification, as they also have a significant impact on object detection models, which are used in applications such as autonomous vehicles and surveillance systems \cite{arxiv:1707.08945}. In these applications, object detection models are used to detect and classify objects in images, and adversarial examples can cause these models to misdetect or misclassify objects. This can have serious consequences, such as in autonomous vehicles, where incorrect object detection can lead to accidents. Furthermore, adversarial examples can be designed to be robust to various transformations, such as rotation and scaling, making them more effective in real-world applications \cite{doi:10.24963/ijcai.2019/134}.

In addition to image classification and object detection, adversarial examples also have an impact on image segmentation models, which are used in applications such as medical imaging and autonomous vehicles \cite{arxiv:1810.05162}. Image segmentation models are used to segment images into different regions, and adversarial examples can cause these models to produce incorrect segmentations. This can have serious consequences, such as in medical imaging, where incorrect segmentations can lead to incorrect diagnoses. The vulnerability of these models to adversarial examples highlights the need for more robust and reliable defense mechanisms.

To mitigate the impact of adversarial examples, researchers have proposed various defense methods, including adversarial training and defensive distillation \cite{doi:10.1145/3133956.3134057}. Adversarial training involves training a model on a dataset that includes adversarial examples, and defensive distillation involves training a model to produce a distilled version of the input data that is more robust to adversarial examples. These defense methods have been shown to be effective in improving the robustness of deep learning models to adversarial examples, but they are not foolproof, and new attack methods are continually being developed to evade them \cite{doi:10.1016/j.eng.2019.12.012}. As the field continues to evolve, it is essential to develop new and more effective defense methods to mitigate the impact of adversarial examples, which will be discussed in the following section.


\subsection{Advanced Topics in Visual Adversarial Attacks}


The realm of visual adversarial attacks has expanded significantly, incorporating various advanced techniques that pose substantial threats to the security and reliability of machine learning models in the physical world. One such technique involves the use of 3D printing to create physical adversarial objects \cite{arxiv:1707.08945}. These objects, when placed in the environment, can deceive machine learning models into misclassifying them, even when viewed from different angles and under various lighting conditions. For instance, a 3D printed stop sign with a subtly altered shape or pattern can be designed to evade detection by autonomous vehicles \cite{arxiv:1807.07769}. The ability to create such objects underscores the importance of considering the physical world implications of adversarial attacks, as highlighted in \cite{doi:10.24963/ijcai.2021/694}.

Another emerging trend in visual adversarial attacks is the utilization of drone-based attacks \cite{doi:10.1109/iros51168.2021.9636638}. By leveraging drones equipped with cameras or projectors, attackers can capture or project images in the physical world, potentially deceiving object detection models used in surveillance or autonomous systems. This approach not only demonstrates the versatility of adversarial attacks but also raises concerns about the security of systems that rely on visual inputs from the environment. 

The concept of universal adversarial perturbations (UAPs) has also been explored in the context of visual adversarial attacks \cite{arxiv:1904.06347}. UAPs refer to a single perturbation that can be applied to any image to cause a deep neural network to misclassify it. The existence of UAPs highlights the vulnerability of deep learning models to adversarial manipulations, even in scenarios where the attacker has limited knowledge of the model's architecture or parameters.

Furthermore, research has shown that adversarial examples can be designed to be robust against various transformations and perturbations, making them more effective in the physical world \cite{doi:10.1109/cvpr.2019.00284}. This includes techniques such as data augmentation and adversarial training, which can enhance the robustness of adversarial examples against different environmental conditions. 

In addition to these techniques, there has been an increasing focus on evaluating the robustness of deep learning models against adversarial attacks in the physical world \cite{doi:10.3390/bdcc8010008}. This includes assessing the effectiveness of different attack methods, such as the Fast Gradient Sign Method (FGSM) and Projected Gradient Descent (PGD), against various deep learning architectures. 

The future of visual adversarial attacks and defenses will likely be shaped by advancements in both fields. As machine learning models become more pervasive in the physical world, the potential impact of adversarial attacks will grow, necessitating the development of more robust and generalizable defense strategies \cite{doi:10.1016/j.eng.2019.12.012}. This may involve the integration of adversarial training with other defense techniques, such as input validation and anomaly detection, to create multi-layered defense systems. Moreover, the increasing use of multimodal sensing and fusion techniques in autonomous systems may offer new avenues for detecting and mitigating adversarial attacks \cite{doi:10.1109/wacv51458.2022.00144}. Ultimately, addressing the challenges posed by visual adversarial attacks will require a concerted effort from researchers, developers, and practitioners to ensure the reliability and security of machine learning systems in the physical world.


\section{Threat Models and Attack Vectors}



\subsection{Introduction to Threat Models}


Threat models play a crucial role in understanding and analyzing the vulnerability of visual adversarial attacks in the physical world. These models simulate various attack scenarios, enabling researchers to evaluate the effectiveness of different attack strategies and develop robust defense mechanisms. In the context of visual adversarial attacks, three primary threat models are commonly employed: white-box, black-box, and gray-box models \cite{doi:10.1145/3128572.3140444}. Each of these models differs in the level of access the attacker has to the target system, influencing the complexity and feasibility of the attack.

The white-box threat model assumes that the attacker has complete knowledge of the target system's architecture, parameters, and training data \cite{arxiv:1707.08945}. This level of access enables the attacker to craft highly effective adversarial examples, as they can precisely calculate the perturbations required to mislead the system. However, this model is often considered unrealistic, as attackers rarely have such extensive knowledge of the target system. Despite this, the white-box model serves as a useful benchmark for evaluating the vulnerability of systems to adversarial attacks.

In contrast, the black-box threat model assumes that the attacker has no knowledge of the target system's internal workings \cite{doi:10.24963/ijcai.2021/694}. Instead, the attacker must rely on indirect methods, such as querying the system with inputs and observing the outputs, to gather information about the system's behavior. This model is more realistic, as attackers often have limited knowledge of the target system. However, black-box attacks can be less effective than white-box attacks, as the attacker must rely on trial and error to craft successful adversarial examples.

The gray-box threat model falls between the white-box and black-box models, assuming that the attacker has some knowledge of the target system, but not complete knowledge \cite{arxiv:1902.06705}. This model is often more realistic than the white-box model, as attackers may have some information about the system's architecture or parameters. Gray-box attacks can be highly effective, as the attacker can use their limited knowledge to craft targeted adversarial examples.

Recent studies have demonstrated the effectiveness of these threat models in evaluating the vulnerability of visual adversarial attacks in the physical world \cite{doi:10.1007/978-3-030-58589-1_3}. For example, \cite{arxiv:1712.09665} showed that a white-box threat model can be used to craft highly effective adversarial patches that can mislead object detectors in the physical world. Similarly, \cite{arxiv:1906.11897} demonstrated the effectiveness of a gray-box threat model in crafting adversarial patches that can evade object detectors.

The development of robust defense mechanisms against visual adversarial attacks requires a thorough understanding of these threat models \cite{arxiv:1810.00069}. By analyzing the strengths and limitations of each model, researchers can develop defense strategies that are effective against a wide range of attack scenarios. For example, \cite{doi:10.1016/j.cose.2020.102092} proposed a defense mechanism that uses a combination of image preprocessing and adversarial training to improve the robustness of image classification systems against adversarial attacks.

In conclusion, threat models play a critical role in understanding and analyzing the vulnerability of visual adversarial attacks in the physical world. By evaluating the effectiveness of different attack strategies and developing robust defense mechanisms, researchers can improve the security and reliability of machine learning systems in the physical world \cite{doi:10.1049/cit2.12028}. Future research should focus on developing more sophisticated threat models that can simulate complex attack scenarios, as well as developing defense mechanisms that can effectively counter these attacks \cite{doi:10.1109/access.2022.3208131}. Additionally, the development of standardized evaluation metrics and methodologies for assessing the effectiveness of visual adversarial attacks and defenses is crucial for advancing the field \cite{arxiv:2010.09670}.


\subsection{Digital Attack Vectors}


Digital attack vectors have emerged as a significant threat to the security of machine learning models, particularly in the context of physical adversarial attacks. As discussed earlier, threat models play a critical role in understanding and analyzing the vulnerability of visual adversarial attacks in the physical world. In the context of digital attack vectors, these threat models can be used to simulate various attack scenarios and evaluate the effectiveness of different attack strategies. For instance, \cite{arxiv:1707.08945} demonstrates the effectiveness of digital attack vectors in generating robust physical-world attacks on deep learning models. The authors propose a general attack algorithm, Robust Physical Perturbations (RP2), to generate robust visual adversarial perturbations under different physical conditions.

Image manipulation is a common digital attack vector used to create adversarial examples \cite{doi:10.24963/ijcai.2021/694}. This can be achieved through various techniques, such as pixel-level perturbations or object insertion \cite{doi:10.1109/access.2018.2807385}. For example, \cite{doi:10.1109/cvprw.2019.00012} shows that adversarial patches can be used to attack person detection systems, highlighting the vulnerability of these systems to digital attack vectors. The use of digital attack vectors, such as image manipulation, can be seen as a precursor to physical attack vectors, which will be discussed in the following section.

Video manipulation is another digital attack vector that has gained significant attention in recent years \cite{doi:10.1109/tcsvt.2023.3341170}. This involves creating adversarial examples by manipulating video frames to deceive machine learning models \cite{doi:10.1007/s44267-024-00066-7}. For instance, \cite{doi:10.1109/tcsvt.2023.3341170} proposes a method to generate adversarial examples for video object segmentation, demonstrating the effectiveness of digital attack vectors in this context. The development of digital attack vectors, including image and video manipulation, has significant implications for the development of robust defense mechanisms against physical adversarial attacks.

The use of generative models, such as Generative Adversarial Networks (GANs), has also been explored as a digital attack vector \cite{doi:10.1109/cvpr46437.2021.01167}. These models can be used to generate realistic adversarial examples that can deceive machine learning models \cite{doi:10.3390/s23020853}. For example, \cite{doi:10.1007/978-3-030-58574-7_41} proposes a method to generate adversarial patches using GANs, highlighting the potential of generative models as a digital attack vector. As the field of digital attack vectors continues to evolve, it is essential to develop effective defense mechanisms that can counter these attacks and protect machine learning models from physical adversarial attacks.

In conclusion, digital attack vectors pose a significant threat to the security of machine learning models, particularly in the context of physical adversarial attacks. The use of image and video manipulation, generative models, and other digital attack vectors can create adversarial examples that can deceive machine learning models \cite{doi:10.1007/s11633-019-1211-x}. As machine learning models become increasingly ubiquitous in various applications, it is essential to develop effective defense mechanisms against digital attack vectors \cite{arxiv:1902.06705}. Future research directions should focus on developing more robust and generalizable defense mechanisms, as well as exploring new digital attack vectors and their potential applications \cite{doi:10.1049/cit2.12028}.


\subsection{Physical Attack Vectors}


Physical attack vectors pose a significant threat to the security of machine learning models in the physical world. These attacks involve manipulating the physical environment to deceive machine learning models, rather than relying on digital manipulations. One of the most common physical attack vectors is the use of adversarial patches, which are physical patches that are designed to be stuck onto an object to mislead a machine learning model \cite{arxiv:1707.07397}. For example, an attacker could print an adversarial patch onto a sticker and stick it onto a stop sign to cause a self-driving car to misclassify the sign \cite{doi:10.1201/9781351251389-8}.

Adversarial patches can be generated using various methods, including optimization-based approaches and generative model-based approaches \cite{doi:10.1145/3128572.3140444}. Optimization-based approaches involve optimizing the patch to maximize the loss function of the machine learning model, while generative model-based approaches involve using a generative model to generate patches that are likely to be misclassified by the model \cite{arxiv:1606.03498}. 

Another physical attack vector is the use of wearable adversarial perturbations, which involve wearing clothing or accessories that are designed to mislead machine learning models \cite{arxiv:1805.06605}. 

Camouflage-based attacks are another type of physical attack vector, which involve using camouflage to hide or disguise an object to mislead a machine learning model \cite{arxiv:1511.05644}. 

Physical attack vectors pose a significant threat to the security of machine learning models in the physical world. These attacks can be used to deceive machine learning models in a variety of applications, including self-driving cars, surveillance systems, and facial recognition systems \cite{doi:10.1201/9781351251389-8}. To defend against these attacks, it is essential to develop robust defense mechanisms that can detect and mitigate physical attack vectors.

One approach to defending against physical attack vectors is to use adversarial training, which involves training a machine learning model on a dataset that includes adversarial examples \cite{arxiv:1901.10513}. Adversarial training can help to improve the robustness of a machine learning model to physical attack vectors, but it requires a large dataset of adversarial examples \cite{arxiv:1902.06705}. 

In conclusion, physical attack vectors pose a significant threat to the security of machine learning models in the physical world. These attacks can be used to deceive machine learning models in a variety of applications, and can be generated using various methods, including optimization-based approaches and generative model-based approaches \cite{arxiv:2001.05574}. To defend against these attacks, it is essential to develop robust defense mechanisms that can detect and mitigate physical attack vectors \cite{doi:10.1109/access.2022.3216291}. 

As \cite{doi:10.1109/eurosp53844.2022.00047} suggests, generating automatic physical examples for machine-learning attacks on computer vision systems is crucial for understanding the vulnerability of these systems. 

Overall, the development of physical attack vectors and defense mechanisms is an ongoing cat-and-mouse game, with each side trying to outmaneuver the other \cite{doi:10.1109/access.2022.3208131}. 

By exploring the current state of physical attack vectors and defense mechanisms, we can better understand the strengths and limitations of these approaches and develop more effective strategies for defending against physical attacks \cite{doi:10.1109/cvpr42600.2020.00090}. 

In the future, we can expect to see the development of more sophisticated physical attack vectors and defense mechanisms, as well as increased awareness of the importance of physical security in machine learning systems \cite{doi:10.1145/3398394}. 

Finally, \cite{doi:10.3390/electronics13132458} highlights the importance of median filtering in enhancing the imperceptibility of perturbations in adversarial examples.


\subsection{Hybrid Attack Vectors}


Hybrid attack vectors combine digital and physical methods to enhance attack effectiveness in the physical world \cite{doi:10.24963/ijcai.2021/694}. This approach allows attackers to create more sophisticated and realistic threats that can evade detection by machine learning models. By leveraging the strengths of both digital and physical methods, hybrid attack vectors can increase the success rate of adversarial attacks and make them more resistant to environmental factors such as lighting and viewpoint changes.

One approach to generating hybrid attack vectors involves using projectors to display adversarial images onto physical objects \cite{doi:10.1109/cvpr42600.2020.01426}. This method enables the creation of dynamic adversarial patterns that can adapt to changing environmental conditions, making them more effective in real-world scenarios. Another approach combines digital perturbations with 3D printing techniques to create physical adversarial objects that can fool machine learning models \cite{arxiv:1707.07397}. These hybrid attack vectors can be designed to evade detection by machine learning models, making them more stealthy and effective \cite{doi:10.1109/cvpr42600.2020.00108}.

The advantages of hybrid attack vectors include increased attack effectiveness and robustness \cite{arxiv:1707.08945}. However, they also have limitations and trade-offs \cite{arxiv:2203.03818}. For example, the use of projectors or 3D printing techniques can increase the complexity and cost of the attack. Additionally, the effectiveness of hybrid attack vectors can depend on the specific machine learning model being targeted, as well as the environmental conditions in which the attack is deployed \cite{arxiv:1902.06705}.

To defend against hybrid attack vectors, it is essential to develop effective defenses, including machine learning-based detection methods and more robust machine learning models \cite{doi:10.1109/tnnls.2021.3105238}. Another approach involves developing more robust machine learning models that are less vulnerable to adversarial attacks \cite{arxiv:1904.13000}. By understanding the properties and implications of hybrid attack vectors, researchers can gain a better understanding of the vulnerabilities of machine learning models and develop more effective defenses against adversarial attacks.

The development of hybrid attack vectors has significant implications for the field of computer vision, particularly in the context of object recognition and detection \cite{doi:10.1007/s10462-024-10841-z}. As machine learning models become increasingly ubiquitous in the physical world, the importance of studying hybrid attack vectors will only continue to grow \cite{doi:10.1109/wacv57701.2024.00393}. By prioritizing the development of more effective defenses against hybrid attack vectors, researchers can help ensure the security and reliability of machine learning models in the physical world \cite{arxiv:2307.15157}. Future research should focus on exploring the properties and implications of hybrid attack vectors, as well as developing more effective defenses against these types of attacks \cite{arxiv:2305.16494}.


\subsection{Evaluation and Comparison of Threat Models and Attack Vectors}


The evaluation and comparison of threat models and attack vectors are crucial steps in understanding the effectiveness and limitations of various adversarial attacks in the physical world. As discussed in \cite{doi:10.1145/3128572.3140444}, the development of robust threat models and attack vectors is essential for improving the security and reliability of deep learning systems. In this subsection, we will evaluate and compare different threat models and attack vectors, including their effectiveness, stealthiness, and robustness in physical adversarial attacks.

Recent studies, such as \cite{arxiv:1707.08945} and \cite{arxiv:1807.07769}, have demonstrated the vulnerability of deep learning models to physical adversarial attacks. These attacks can be launched using various threat models, including white-box, black-box, and gray-box models. The white-box model assumes that the attacker has full access to the model's parameters and architecture, while the black-box model assumes that the attacker has no knowledge of the model's internal workings. The gray-box model falls between these two extremes, where the attacker has some knowledge of the model's parameters and architecture.

In terms of attack vectors, \cite{doi:10.1109/cvpr.2019.00284} and \cite{doi:10.24963/ijcai.2019/134} have shown that input diversity and transferability are essential factors in improving the effectiveness of adversarial attacks. The use of diverse input patterns and transferable attack vectors can increase the success rate of adversarial attacks, making them more challenging to detect and defend against.

The evaluation of threat models and attack vectors can be performed using various metrics, including success rate, detection rate, and robustness metrics. As discussed in \cite{doi:10.3390/bdcc8010008}, the success rate of an attack is a critical metric in evaluating its effectiveness. The detection rate, on the other hand, measures the ability of a defense mechanism to detect and prevent adversarial attacks. Robustness metrics, such as adversarial accuracy and certified robustness, can provide a more comprehensive understanding of a model's robustness to adversarial attacks.

Comparing different threat models and attack vectors, \cite{doi:10.1016/j.eng.2019.12.012} and \cite{doi:10.3390/electronics10172132} have shown that the choice of threat model and attack vector can significantly impact the effectiveness of an adversarial attack. For example, white-box models can be more effective than black-box models, but they require more knowledge of the model's internal workings. Similarly, transferable attack vectors can be more effective than non-transferable vectors, but they require more computational resources.

In conclusion, the evaluation and comparison of threat models and attack vectors are essential steps in understanding the effectiveness and limitations of various adversarial attacks in the physical world. By analyzing the strengths and weaknesses of different threat models and attack vectors, researchers can develop more robust defense mechanisms and improve the security and reliability of deep learning systems. As discussed in \cite{doi:10.1145/3702638}, future research directions include the development of more robust threat models and attack vectors, as well as the evaluation of their effectiveness in various applications, including medical image analysis and autonomous driving. By addressing these challenges, researchers can contribute to the development of more secure and reliable deep learning systems that can withstand the threats of adversarial attacks.

Furthermore, \cite{doi:10.24963/ijcai.2021/694} and \cite{doi:10.1016/j.sysarc.2020.101766} have highlighted the importance of considering the physical world constraints when evaluating threat models and attack vectors. The physical world constraints, such as lighting and viewpoint changes, can significantly impact the effectiveness of adversarial attacks. Therefore, it is essential to consider these constraints when developing and evaluating threat models and attack vectors.

In addition, \cite{arxiv:1803.00404} and \cite{doi:10.1016/j.eng.2019.12.012} have shown that adversarial training can be an effective method for improving the robustness of deep learning models. Adversarial training involves training the model on a dataset that includes adversarial examples, which can help the model learn to recognize and defend against these examples. By combining adversarial training with other defense mechanisms, researchers can develop more robust deep learning systems that can withstand the threats of adversarial attacks.

Overall, the evaluation and comparison of threat models and attack vectors are critical steps in understanding the effectiveness and limitations of various adversarial attacks in the physical world. By analyzing the strengths and weaknesses of different threat models and attack vectors, researchers can develop more robust defense mechanisms and improve the security and reliability of deep learning systems. As discussed in \cite{doi:10.1007/s11263-023-01934-3}, future research directions include the development of more robust threat models and attack vectors, as well as the evaluation of their effectiveness in various applications, including medical image analysis and autonomous driving.


\subsection{Future Directions and Open Challenges}


The development of threat models and attack vectors for physical adversarial attacks has been a rapidly evolving field, with significant advancements in recent years. As researchers continue to explore new methods for generating adversarial examples, it is essential to identify future directions and open challenges in this area. Building on the evaluation and comparison of threat models and attack vectors discussed in the previous subsection, one of the primary challenges is the development of more robust and generalizable defense strategies that can effectively counter physical adversarial attacks \cite{doi:10.1109/access.2023.3345000}. Current defense mechanisms often focus on detecting and mitigating specific types of attacks, but they may not be effective against more sophisticated or unknown attacks \cite{doi:10.24963/ijcai.2021/694}.

To address this challenge, researchers have proposed various approaches, including the use of adversarial training \cite{arxiv:2210.07540}, defensive distillation \cite{doi:10.1109/cvprw56347.2022.00034}, and certification methods \cite{arxiv:1707.08945}. However, these approaches have their limitations, and there is a need for more innovative and effective defense strategies. For example, the use of generative models, such as Generative Adversarial Networks (GANs) and Variational Autoencoders (VAEs), has shown promise in generating robust adversarial examples \cite{doi:10.1109/cvpr.2019.00935}. Nevertheless, the development of more advanced defense mechanisms that can effectively counter these attacks is an open challenge \cite{doi:10.58496/adsa/2024/011}.

As the field continues to evolve, another future direction is the exploration of new threat models and attack vectors, such as the use of 3D printing and drone-based attacks \cite{arxiv:1907.05418}. These attacks have the potential to be more effective and stealthy than traditional attacks, and they pose significant challenges to the development of robust defense mechanisms. Additionally, the use of adversarial examples in the physical world raises concerns about the safety and security of autonomous systems, such as self-driving cars and drones \cite{doi:10.1109/access.2023.3321118}. Therefore, it is essential to develop more effective defense strategies that can counter these types of attacks, which will be crucial in ensuring the reliability and security of deep learning systems in various applications, including those discussed in the following subsection.

The development of more robust and generalizable defense mechanisms also requires a better understanding of the properties of adversarial examples \cite{doi:10.1145/3274895.3274904}. For example, the use of metrics such as peak signal-to-noise ratio (PSNR) and structural similarity index (SSIM) can help evaluate the similarity between adversarial and original examples. However, there is a need for more advanced metrics that can effectively capture the properties of adversarial examples in the physical world. Furthermore, the development of more effective defense mechanisms requires a better understanding of the attack strategies and methodologies used by adversaries \cite{arxiv:1902.10899}. By addressing these challenges and developing more effective defense mechanisms, researchers can contribute to the development of more secure and reliable deep learning systems that can withstand the threats of physical adversarial attacks.

In conclusion, the development of threat models and attack vectors for physical adversarial attacks is a rapidly evolving field, with significant advancements in recent years. However, there are still many open challenges and future directions that need to be addressed, including the development of more robust and generalizable defense strategies, the exploration of new threat models and attack vectors, and the development of more advanced metrics for evaluating adversarial examples. By addressing these challenges, researchers can develop more effective defense mechanisms that can counter physical adversarial attacks and ensure the safety and security of autonomous systems \cite{doi:10.1109/tpami.2024.3430860}. As noted in \cite{doi:10.1016/j.eswa.2025.128967}, the development of more effective defense mechanisms requires a multidisciplinary approach that combines advances in computer vision, machine learning, and cybersecurity. Ultimately, the goal is to develop more robust and secure autonomous systems that can operate effectively in the presence of physical adversarial attacks \cite{arxiv:2409.20426}.


\section{Defense Mechanisms and Strategies}



\subsection{Detection-Based Defenses}


Detection-based defenses are a crucial component in the arsenal against visual adversarial attacks in the physical world. These defenses aim to identify and flag potential adversarial examples before they can cause harm to the system. The scope of this subsection encompasses various detection methods, including statistical approaches, machine learning-based techniques, and signal processing strategies. By exploring these different methodologies, researchers and practitioners can better understand the strengths and limitations of each approach and develop more effective countermeasures against visual adversarial attacks.

Statistical methods for detecting adversarial examples have been extensively studied, with techniques such as hypothesis testing and anomaly detection being employed to identify outliers in the data distribution \cite{doi:10.1145/3128572.3140444}. These methods are often based on the assumption that adversarial examples will have different statistical properties than legitimate inputs. However, as shown in \cite{arxiv:1902.06705}, the effectiveness of these methods can be limited by the complexity of the data distribution and the sophistication of the adversarial attacks.

Machine learning-based approaches have also been widely adopted for detecting adversarial examples. These methods typically involve training a separate classifier to distinguish between legitimate and adversarial inputs \cite{arxiv:1810.00069}. The key advantage of these approaches is their ability to learn complex patterns in the data that may not be apparent through statistical analysis alone. However, as noted in \cite{doi:10.1016/j.cose.2020.102092}, the performance of these methods can be highly dependent on the quality of the training data and the specific attack scenario.

Signal processing techniques have also been explored for detecting adversarial examples, with methods such as filtering and transformation being used to remove or mitigate the effects of adversarial perturbations \cite{arxiv:1707.08945}. These techniques can be particularly effective in scenarios where the adversarial perturbations are relatively simple or have a distinct frequency signature. However, as shown in \cite{doi:10.1109/tnnls.2023.3321432}, more sophisticated attacks can be designed to evade these detection methods.

In addition to these specific techniques, several papers have proposed more general frameworks for detecting adversarial examples. For example, \cite{arxiv:2005.10884} presents a framework for detecting adversarial patches using a combination of statistical and machine learning-based approaches. Similarly, \cite{doi:10.1109/cvpr52688.2022.01455} proposes a framework for detecting adversarial patches in object detection scenarios.

Despite the progress made in detecting adversarial examples, several challenges and limitations remain. One of the primary challenges is the need for a robust and generalizable detection method that can effectively identify adversarial examples across a wide range of attack scenarios and datasets. As noted in \cite{doi:10.1109/access.2022.3216291}, the development of such a method is an active area of research, with several promising approaches being explored.

Another challenge is the potential for adversarial examples to be designed specifically to evade detection. As shown in \cite{doi:10.1007/s11633-019-1211-x}, adversarial attacks can be crafted to have a minimal impact on the input data while still causing misclassification. This highlights the need for detection methods that can effectively identify subtle changes in the input data.

In conclusion, detection-based defenses are a critical component in the defense against visual adversarial attacks in the physical world. While various statistical, machine learning-based, and signal processing techniques have been proposed, several challenges and limitations remain. Further research is needed to develop more effective and generalizable detection methods that can identify adversarial examples across a wide range of attack scenarios and datasets. As noted in \cite{doi:10.1109/access.2024.3435347}, the development of such methods is essential for ensuring the security and reliability of deep learning systems in the physical world. By continuing to explore and develop new detection methods, researchers and practitioners can help to mitigate the risks associated with visual adversarial attacks and ensure the safe deployment of deep learning systems in a wide range of applications \cite{doi:10.1016/j.eswa.2025.128967}.


\subsection{Mitigation Strategies}


Mitigation strategies against visual adversarial attacks are crucial for ensuring the reliability and security of machine learning models in the physical world. As discussed in the previous section, detection-based defenses are a critical component in the defense against visual adversarial attacks, and mitigation strategies can be seen as a complementary approach to enhance the security of machine learning models. These strategies can be broadly categorized into input preprocessing, data augmentation, and robust optimization methods. Input preprocessing techniques, such as data normalization and feature extraction, can help mitigate the effects of adversarial attacks by reducing the impact of perturbations on the input data \cite{arxiv:1712.09665}. Data augmentation methods, including randomization and perturbation techniques, can also improve the robustness of machine learning models against visual adversarial attacks \cite{doi:10.24963/ijcai.2021/694}.

Robust optimization methods, such as adversarial training and robust loss functions, have been shown to be effective in improving the security of machine learning models in the physical world \cite{arxiv:1707.08945}. Adversarial training involves training the model on a dataset that includes adversarial examples, which can help improve the model's robustness to attacks \cite{doi:10.1016/j.eng.2019.12.012}. Robust loss functions, such as the hinge loss and the triplet loss, can also be used to improve the model's robustness to attacks \cite{arxiv:1902.06705}. These methods can be used in conjunction with detection-based defenses to provide a multi-layered defense against visual adversarial attacks.

In addition to these methods, several other mitigation strategies have been proposed, including the use of generative models to detect and remove adversarial perturbations \cite{doi:10.1109/cvpr46437.2021.01167}, and the use of attention mechanisms to focus on the most relevant features of the input data \cite{doi:10.1109/cvpr52688.2022.01480}. These methods have been shown to be effective in mitigating the effects of visual adversarial attacks, but they often require significant computational resources and may not be feasible for real-time applications. Recent studies have also explored the use of physical-world constraints to improve the robustness of machine learning models to adversarial attacks \cite{doi:10.1109/access.2023.3321118}. For example, the use of physical-world constraints, such as the laws of physics and the properties of materials, can help improve the robustness of machine learning models to attacks \cite{doi:10.1109/cvpr52729.2023.01187}.

The development of mitigation strategies against visual adversarial attacks is closely related to the development of robustness enhancement techniques, which will be discussed in the following section. Robustness enhancement techniques, such as adversarial training and defensive distillation, can be used to improve the model's robustness to adversarial attacks, and can be seen as a natural extension of the mitigation strategies discussed in this section. Despite the advances in mitigation strategies, developing effective methods that can defend against multiple types of attacks, such as both digital and physical attacks, is an important area of research \cite{doi:10.1049/cit2.12028}. Additionally, developing mitigation strategies that can be used in real-time applications, such as autonomous vehicles and surveillance systems, is crucial for ensuring the safety and security of these systems \cite{doi:10.1109/access.2024.3418614}.

In conclusion, mitigation strategies against visual adversarial attacks are essential for ensuring the reliability and security of machine learning models in the physical world. Several methods, including input preprocessing, data augmentation, and robust optimization methods, have been proposed to mitigate the effects of these attacks. However, developing effective mitigation strategies that can defend against multiple types of attacks and can be used in real-time applications remains an important area of research \cite{doi:10.1007/s10462-024-10841-z}. Future studies should focus on developing more robust and efficient mitigation strategies that can be used in a variety of applications, including autonomous vehicles, surveillance systems, and other safety-critical systems \cite{doi:10.1038/s41598-025-15753-8}.


\subsection{Robustness Enhancement Techniques}


Robustness enhancement techniques are crucial for improving the security of machine learning models in the physical world. These techniques aim to increase the model's resistance to adversarial attacks, which are designed to mislead the model by adding subtle perturbations to the input data. One of the most effective robustness enhancement techniques is adversarial training \cite{doi:10.1145/3128572.3140444}, which involves training the model on a mixture of clean and adversarial examples. This approach has been shown to improve the model's robustness to various types of attacks, including those that are designed to be transferable across different models \cite{doi:10.1109/cvpr.2019.00284}.

Another technique that has gained significant attention in recent years is defensive distillation \cite{arxiv:1805.06605}, which involves training a model to produce a condensed representation of the input data that is more robust to adversarial attacks. This approach has been shown to be effective in improving the robustness of models to various types of attacks, including those that are designed to be transferable across different models \cite{doi:10.1609/aaai.v32i1.11828}. Certification methods, such as \cite{arxiv:1902.06705}, have also been proposed to provide a formal guarantee of the model's robustness to adversarial attacks.

In addition to these techniques, several other approaches have been proposed to enhance the robustness of machine learning models, including input preprocessing \cite{arxiv:1711.00117}, data augmentation \cite{doi:10.1109/cvpr.2019.00284}, and robust optimization methods \cite{arxiv:1911.05268}. These approaches have been shown to be effective in improving the robustness of models to various types of attacks, including those that are designed to be transferable across different models \cite{doi:10.1016/j.eng.2019.12.012}.

Despite the progress made in developing robustness enhancement techniques, there are still several challenges that need to be addressed. One of the main challenges is the trade-off between robustness and accuracy, as increasing the model's robustness to adversarial attacks can often come at the cost of reducing its accuracy on clean data \cite{doi:10.1109/cvpr42600.2020.00090}. Another challenge is the lack of standardization in evaluating the robustness of models to adversarial attacks, which can make it difficult to compare the effectiveness of different techniques \cite{doi:10.3390/bdcc8010008}.

To address these challenges, researchers have proposed several emerging trends and techniques, including the use of generative models \cite{arxiv:2001.05574} and transfer learning \cite{doi:10.1109/access.2022.3216291} to improve the robustness of models to adversarial attacks. These approaches have shown promising results in improving the robustness of models to various types of attacks, including those that are designed to be transferable across different models \cite{doi:10.1109/eurosp53844.2022.00047}.

In conclusion, robustness enhancement techniques are essential for improving the security of machine learning models in the physical world. While significant progress has been made in developing these techniques, there are still several challenges that need to be addressed. Emerging trends and techniques, such as the use of generative models and transfer learning, have shown promising results in improving the robustness of models to adversarial attacks. Further research is needed to develop more effective and efficient robustness enhancement techniques that can provide a formal guarantee of the model's robustness to adversarial attacks \cite{doi:10.1109/sp46215.2023.10179316}. By synthesizing the insights from \cite{doi:10.1109/access.2022.3208131} and \cite{doi:10.1007/s10462-024-10841-z}, we can develop a more comprehensive understanding of the challenges and opportunities in enhancing the robustness of machine learning models to adversarial attacks.


\subsection{Physical World Considerations}


The physical world poses unique challenges for visual adversarial attacks, as environmental factors and the need for realism can significantly impact the effectiveness of such attacks. As noted in \cite{doi:10.24963/ijcai.2021/694}, adversarial examples that are successful in the digital realm may not necessarily translate to the physical world, where factors such as lighting, weather, and viewpoint changes can affect the attack's success. Furthermore, \cite{arxiv:1707.08945} highlights the importance of considering the physical world's constraints when generating adversarial examples, as attacks that are robust in the digital realm may not be as effective in the physical world. This emphasizes the need for robustness enhancement techniques, such as those discussed in the previous section, to be adapted to the physical world context.

One key consideration in physical world adversarial attacks is the need for realism. As \cite{doi:10.1109/cvpr42600.2020.00108} notes, adversarial examples that are not realistic may be easily detectable by humans, which can limit their effectiveness. To address this, researchers have proposed various methods for generating more realistic adversarial examples, such as \cite{doi:10.1109/cvpr42600.2020.01426}, which uses a generative model to create adversarial examples that are more resilient to physical world transformations. This is closely related to the development of robustness enhancement techniques, such as defensive distillation and certification methods, which can provide a formal guarantee of the model's robustness to adversarial attacks.

Another important factor in physical world adversarial attacks is the impact of environmental factors. As \cite{arxiv:2203.03818} demonstrates, even natural phenomena such as shadows can be used to create effective adversarial attacks. This highlights the need for defenses that can account for a wide range of environmental factors, as noted in \cite{arxiv:2408.01541}. The evaluation of such defenses is crucial, and will be discussed in the following section, which will explore the challenges and opportunities in evaluating the effectiveness of visual adversarial attacks and defenses in the physical world.

In addition to these challenges, physical world adversarial attacks also raise important questions about the evaluation of such attacks. As \cite{doi:10.1109/tpami.2025.3610085} notes, existing evaluation metrics may not be sufficient for assessing the effectiveness of physical world adversarial attacks, and new metrics that account for physical world constraints may be needed. Furthermore, \cite{doi:10.1109/sp46215.2023.10179316} highlights the importance of considering the trade-offs between different evaluation metrics, as well as the need for more comprehensive evaluation frameworks. This emphasizes the need for continued innovation in evaluation metrics and methodologies to keep pace with the evolving landscape of visual adversarial attacks and defenses.

Despite these challenges, physical world adversarial attacks also offer opportunities for innovation and advancement in the field of computer vision. As \cite{arxiv:2012.12235} notes, the study of physical world adversarial attacks can provide valuable insights into the design of more robust and secure computer vision systems. Furthermore, \cite{doi:10.1609/aaai.v35i9.16927} demonstrates the potential for adversarial training to improve the robustness of computer vision models to natural perturbations, which can be particularly important in physical world scenarios. This highlights the potential for physical world adversarial attacks to inform the development of more secure and robust systems, and emphasizes the need for further research in this area.

In conclusion, physical world adversarial attacks pose unique challenges and opportunities for the field of computer vision. By considering the impact of environmental factors and the need for realism, researchers can develop more effective and robust defenses against such attacks. As noted in \cite{arxiv:1901.10513}, the study of physical world adversarial attacks can also provide valuable insights into the fundamental limitations of computer vision models, and can inform the development of more secure and robust systems. Future research directions may include the development of more comprehensive evaluation frameworks, as well as the exploration of new methods for generating realistic and effective adversarial examples in the physical world, such as those proposed in \cite{arxiv:2305.16494} and \cite{doi:10.3389/fnbot.2023.1129720}.


\subsection{Evaluation Metrics and Methodologies}


Evaluating the effectiveness of visual adversarial attacks and defenses in the physical world is a complex task that requires careful consideration of various metrics and methodologies. As discussed in \cite{doi:10.1145/3128572.3140444}, the development of robust evaluation metrics is crucial for assessing the success of adversarial attacks and the effectiveness of defense mechanisms. In the context of physical-world attacks, metrics such as misclassification rates, detection evasion rates, and robustness metrics (e.g., adversarial accuracy and certified robustness) are essential for evaluating the performance of visual adversarial attacks and defenses \cite{arxiv:1707.08945}.

One of the primary challenges in evaluating visual adversarial attacks and defenses is the lack of standardized metrics and methodologies. Different studies often employ varying evaluation protocols, making it difficult to compare results and draw conclusive insights \cite{doi:10.1109/access.2018.2807385}. To address this issue, researchers have proposed several evaluation frameworks, including the use of digital and physical testing methodologies \cite{doi:10.3390/bdcc8010008}. These frameworks provide a comprehensive approach to evaluating the effectiveness of visual adversarial attacks and defenses, considering factors such as attack success rates, defense robustness, and environmental conditions.

In addition to evaluation metrics and methodologies, the choice of attack and defense strategies is critical in assessing the effectiveness of visual adversarial attacks and defenses. As shown in \cite{doi:10.1109/cvpr.2019.00284}, input diversity can significantly improve the transferability of adversarial examples, making them more effective in fooling deep learning models. Similarly, defense strategies such as adversarial training, defensive distillation, and certification methods can enhance the robustness of deep learning models against adversarial attacks \cite{arxiv:1803.00404}.

Recent studies have also explored the use of emerging trends and techniques, such as 3D printing and drone-based attacks, to create more sophisticated and realistic visual adversarial attacks \cite{arxiv:1907.05418}. These advances highlight the need for continued innovation in evaluation metrics and methodologies to keep pace with the evolving landscape of visual adversarial attacks and defenses. Furthermore, the development of more robust and generalizable defense mechanisms is essential for ensuring the security and reliability of deep learning models in the physical world \cite{doi:10.3390/electronics10172132}.

The evaluation of visual adversarial attacks and defenses is not limited to the technical aspects of deep learning models. As discussed in \cite{doi:10.3390/cancers15174228}, the impact of adversarial attacks on medical image analysis can have significant consequences for patient care and safety. Therefore, it is essential to consider the broader implications of visual adversarial attacks and defenses, including their potential effects on human life and safety \cite{doi:10.1145/3702638}.

In conclusion, the evaluation of visual adversarial attacks and defenses in the physical world is a complex and multifaceted task that requires careful consideration of various metrics, methodologies, and strategies. By synthesizing insights from existing studies, such as \cite{doi:10.24963/ijcai.2021/694} and \cite{arxiv:1807.07769}, researchers can develop more effective evaluation frameworks and defense mechanisms to enhance the security and reliability of deep learning models in the physical world. As the field continues to evolve, it is essential to prioritize innovation, standardization, and collaboration to address the challenges and opportunities presented by visual adversarial attacks and defenses \cite{doi:10.1007/s11263-023-01934-3}.


\subsection{Open Challenges and Future Directions}


The field of visual adversarial attacks and defenses in the physical world is rapidly evolving, with new challenges and opportunities emerging continuously. As the previous discussions have highlighted the importance of evaluating the effectiveness of visual adversarial attacks and defenses, it is essential to address the open challenges and future directions that remain to be explored. One of the primary challenges is the development of more robust and generalizable defense mechanisms that can effectively counter physical adversarial attacks in various environments and conditions \cite{doi:10.1007/s11633-019-1211-x}. Current defenses often rely on digital testing methodologies, which may not accurately capture the complexities of real-world scenarios \cite{doi:10.1007/978-3-030-58589-1_3}. Therefore, there is a need for more comprehensive evaluation methodologies that incorporate physical testing and consider factors such as lighting, weather, and viewpoint changes.

Building on the need for more effective evaluation methodologies, another significant challenge is the lack of standardized metrics for evaluating the effectiveness of physical adversarial attacks and defenses \cite{doi:10.1109/tgrs.2022.3225306}. Most existing metrics focus on digital attacks and may not be directly applicable to physical-world scenarios. The development of robust and relevant metrics is crucial for comparing different attack and defense strategies and assessing their performance in real-world settings. Furthermore, the exploration of new threat models and attack vectors, such as 3D printing and drone-based attacks, is essential for understanding the potential risks and vulnerabilities of physical adversarial attacks \cite{arxiv:1707.08945}. This, in turn, can inform the development of more effective defense mechanisms that can mitigate these threats.

In addition to these technical challenges, the use of generative models, such as Generative Adversarial Networks (GANs) and Variational Autoencoders (VAEs), has shown promise in generating realistic adversarial examples \cite{doi:10.1109/cvpr.2019.00935}. However, these models often require significant computational resources and may not be feasible for real-time applications \cite{arxiv:1907.05418}. The development of more efficient and effective generative models is necessary for generating high-quality adversarial examples that can be used in various physical-world scenarios \cite{doi:10.1016/j.jisa.2020.102694}. Moreover, there are also ethical and regulatory considerations that need to be addressed, as the development of physical adversarial attacks and defenses raises concerns about privacy, security, and safety, particularly in applications such as autonomous vehicles and surveillance systems \cite{arxiv:2007.16118}.

As the field continues to evolve, future research directions include the exploration of new application domains, such as robotics and healthcare, and the development of more sophisticated attack and defense strategies \cite{doi:10.1016/j.eng.2019.12.012}. The use of transfer learning and meta-learning techniques may also be beneficial for improving the robustness and generalizability of physical adversarial defenses \cite{arxiv:2210.07540}. Furthermore, the development of more realistic and effective physical adversarial attacks, such as those using 3D printing and projector-based attacks, is essential for evaluating the robustness of physical adversarial defenses \cite{doi:10.1109/cvpr52688.2022.01487}. By addressing these challenges and exploring new research directions, we can develop more robust and generalizable defense mechanisms and ensure the safe and reliable deployment of machine learning models in various physical-world applications \cite{doi:10.1109/tpami.2024.3430860}. Ultimately, synthesizing information from various studies, including \cite{doi:10.1201/9781351251389-8}, \cite{doi:10.1609/aaai.v34i01.5443}, and \cite{doi:10.1109/tpami.2022.3176760}, will be crucial for gaining a deeper understanding of the current state of the field and identifying potential avenues for future research and development.


\section{Physical World Considerations and Constraints}



\subsection{Environmental Factors in Physical Adversarial Attacks}


The success of visual adversarial attacks in the physical world is significantly influenced by various environmental factors, including lighting, weather, and viewpoint changes. These factors can affect the visibility and effectiveness of adversarial patches or attacks, making it challenging to design robust defenses. As noted in \cite{doi:10.1145/3128572.3140444}, environmental factors can be exploited to create more effective adversarial attacks. For instance, \cite{arxiv:1707.08945} demonstrates that adversarial examples can be designed to be robust against different lighting conditions, making them more effective in real-world scenarios.

Lighting conditions, in particular, play a crucial role in the success of visual adversarial attacks. Different lighting conditions, such as daylight, nighttime, indoor, or outdoor lighting, can affect the visibility of adversarial patches or attacks. As shown in \cite{doi:10.24963/ijcai.2021/694}, adversarial examples can be designed to be more effective in certain lighting conditions, making them more challenging to detect. Furthermore, \cite{arxiv:1906.11897} highlights the importance of considering lighting conditions when designing adversarial patches, as they can significantly impact the attack's effectiveness.

Weather conditions, such as rain, fog, or snow, can also affect the success of visual adversarial attacks. As demonstrated in \cite{doi:10.1109/tnnls.2023.3321432}, weather conditions can reduce the effectiveness of adversarial attacks, making them less robust in real-world scenarios. However, \cite{doi:10.1016/j.eswa.2025.128967} shows that adversarial patches can be designed to be robust against certain weather conditions, such as rain or fog, making them more effective in real-world scenarios.

Viewpoint changes, such as changes in the orientation or position of the camera, can also impact the success of visual adversarial attacks. As noted in \cite{arxiv:1902.06705}, viewpoint changes can affect the effectiveness of adversarial attacks, making them less robust in real-world scenarios. However, \cite{doi:10.1007/978-3-030-10925-7_4} demonstrates that adversarial attacks can be designed to be robust against viewpoint changes, making them more effective in real-world scenarios.

In addition to these environmental factors, the material and texture of adversarial objects or patches can also impact their effectiveness. As shown in \cite{doi:10.1007/978-3-030-58574-7_41}, the texture and material of adversarial patches can affect their visibility and effectiveness, making them more challenging to detect. Furthermore, \cite{doi:10.1007/978-3-030-58589-1_3} highlights the importance of considering the material and texture of adversarial objects or patches when designing defenses, as they can significantly impact the attack's effectiveness.

To design effective defenses against visual adversarial attacks, it is essential to consider these environmental factors and their impact on the attack's effectiveness. As noted in \cite{doi:10.1109/access.2022.3216291}, defenses can be designed to be more robust against adversarial attacks by considering the environmental factors that affect their effectiveness. Furthermore, \cite{doi:10.1007/s11633-019-1211-x} highlights the importance of developing more robust and generalizable defenses that can effectively detect and mitigate adversarial attacks in various environmental conditions.

In conclusion, environmental factors, such as lighting, weather, and viewpoint changes, play a crucial role in the success of visual adversarial attacks in the physical world. To design effective defenses, it is essential to consider these factors and their impact on the attack's effectiveness. By developing more robust and generalizable defenses, we can improve the security and reliability of machine learning models in the physical world, making them more resilient to adversarial attacks. As future research directions, it is essential to explore more advanced defense mechanisms that can effectively detect and mitigate adversarial attacks in various environmental conditions, such as \cite{doi:10.1109/access.2022.3208131} and \cite{doi:10.1109/tnnls.2025.3561225}.


\subsection{Realism and Stealth in Physical Adversarial Attacks}


The design of effective physical adversarial attacks relies on two crucial factors: realism and stealth. Realism refers to the ability of the adversarial patch or object to blend seamlessly into its surroundings, making it difficult for humans to detect. Stealth, on the other hand, pertains to the patch's or object's ability to evade detection by automated systems. Both factors are essential in ensuring the success of physical adversarial attacks, as they enable the attacker to manipulate the target system without arousing suspicion \cite{arxiv:1712.09665}. 

As discussed earlier, environmental factors such as lighting, weather, and viewpoint changes can significantly impact the effectiveness of adversarial attacks. To achieve realism and stealth, attackers must consider these factors when designing their attacks. For instance, \cite{arxiv:1902.06705} highlights the importance of evaluating adversarial robustness under various environmental conditions. By taking these factors into account, attackers can design more effective physical adversarial attacks. 

One approach to achieving realism and stealth is through the use of 3D modeling and rendering techniques \cite{doi:10.1109/cvpr52688.2022.01487}. By creating highly realistic models of the adversarial object or patch, attackers can design attacks that are more likely to succeed in the physical world. For example, \cite{arxiv:1707.08945} demonstrates the effectiveness of using 3D printed objects to create physical adversarial examples that can fool deep learning models. Similarly, \cite{arxiv:1806.02299} shows that using textured patches can improve the stealthiness of physical adversarial attacks.

Recent studies have also explored the use of generative models to improve the realism and stealth of physical adversarial attacks \cite{doi:10.1109/cvpr46437.2021.01167}. By leveraging the capabilities of generative models, attackers can create highly realistic and stealthy adversarial patches that are more likely to succeed in the physical world. For instance, \cite{doi:10.1007/978-3-030-58574-7_41} demonstrates the effectiveness of using generative models to create textured patches that can fool object detectors.

The development of more sophisticated attack methods, such as the use of projectors to display adversarial images \cite{doi:10.1109/vr51125.2022.00073}, has highlighted the need for more robust defense mechanisms. The use of projectors allows attackers to create dynamic and stealthy adversarial attacks that can be more difficult to detect. Additionally, \cite{doi:10.1109/tpami.2022.3176760} proposes a novel approach to creating stealthy physical adversarial attacks using stickers. The method involves manipulating the sticker's pasting position and rotation angle on the target object to perform physical attacks.

The importance of realism and stealth in physical adversarial attacks is also highlighted in \cite{arxiv:2007.16118}. The study demonstrates the effectiveness of using physical adversarial attacks to deceive vehicle detectors in a simulated environment. The results show that the attacks can be highly effective, even when the adversarial patch is small and stealthy. As the field continues to evolve, it is essential to develop more robust defense mechanisms that can detect and mitigate these types of attacks \cite{arxiv:1810.00069}.

In conclusion, the design of effective physical adversarial attacks requires careful consideration of realism and stealth. By leveraging techniques such as 3D modeling, generative models, and projectors, attackers can create highly realistic and stealthy attacks that can fool deep learning models in the physical world. The pursuit of realism in physical adversarial attacks will be further discussed in the following section, which explores the techniques for enhancing realism in physical adversarial attacks \cite{arxiv:1707.07397} and \cite{doi:10.1016/j.neucom.2022.05.031}. Ultimately, the goal of this research is to improve the security and robustness of deep learning models in the physical world, and to develop more effective defense mechanisms against physical adversarial attacks \cite{doi:10.1109/tnnls.2025.3561225}.


\subsection{Techniques for Enhancing Realism in Physical Adversarial Attacks}


The pursuit of realism in physical adversarial attacks has become a critical aspect of research in the field of visual adversarial attacks and defenses. As \cite{arxiv:1707.07397} and \cite{doi:10.1016/j.neucom.2022.05.031} have demonstrated, the use of 3D modeling and rendering can significantly enhance the realism of physical adversarial attacks. By creating a 3D digital approximation of the scene, attackers can improve the adversarial patch's impact in real-world settings. Furthermore, \cite{arxiv:2203.09123} highlights the importance of object-based diverse input in generating more realistic and transferable adversarial examples.

Another approach to enhancing realism is through the use of projectors or other technologies to create dynamic adversarial patterns. As \cite{arxiv:2001.05574} and \cite{doi:10.1109/eurosp53844.2022.00047} have shown, projectors can be used to dynamically generate adversarial patterns on objects or surfaces, making it more challenging for defenders to detect and mitigate the attacks. Moreover, \cite{doi:10.1609/aaai.v34i07.6918} demonstrates the effectiveness of using super-pixels to decrease the number of pixels that need to compute gradients, thereby improving the efficiency of adversarial attacks on video recognition models.

The use of evolutionary algorithms, such as those proposed in \cite{doi:10.1186/s41074-019-0053-3}, can also enhance the realism of physical adversarial attacks. By exploring the latent space of a generative adversarial network with an evolutionary algorithm, attackers can find vectors representing adversarial attacks that are more realistic and effective. Additionally, \cite{doi:10.1109/access.2024.3415356} highlights the potential of using gradient-based and Bayesian-optimized approaches to generate adversarial examples that are more effective and realistic.

In terms of emerging trends and challenges, \cite{doi:10.1109/sp46215.2023.10179316} raises important questions about the effectiveness of unrealistic adversarial hardening against realistic adversarial attacks. The study suggests that unrealistic adversarial examples can be used to protect models against realistic examples, but the results are inconsistent across different use cases. Furthermore, \cite{doi:10.1609/aaai.v38i19.30150} demonstrates the potential of using visual adversarial examples to jailbreak aligned large language models, highlighting the need for more robust defenses against adversarial attacks.

As research in this area continues to evolve, it is essential to consider the practical implications of these techniques. \cite{arxiv:1805.06605} and \cite{doi:10.1109/access.2022.3160283} propose the use of generative adversarial networks to defend against adversarial attacks, highlighting the potential for more robust and realistic defenses. Moreover, \cite{doi:10.1109/access.2022.3216291} provides a comprehensive overview of the existing methods for adversarial robustness, highlighting the need for more efficient and effective approaches to defending against adversarial attacks.

In conclusion, the techniques for enhancing realism in physical adversarial attacks are diverse and rapidly evolving. By leveraging 3D modeling, rendering, projectors, and evolutionary algorithms, attackers can create more realistic and effective adversarial examples. However, defenders must also adapt and develop more robust defenses against these attacks. As \cite{doi:10.1109/access.2022.3208131} and \cite{doi:10.1007/s10462-024-10841-z} highlight, the development of more efficient and effective methods for adversarial robustness is crucial to ensuring the security and reliability of machine learning models in the physical world. Future research should focus on developing more realistic and effective adversarial attacks, as well as more robust and efficient defenses against these attacks.


\subsection{Evaluation Metrics for Physical Adversarial Attacks}


The development of robust evaluation metrics for physical adversarial attacks is crucial for assessing the effectiveness and realism of these attacks, as highlighted in the previous discussion on enhancing realism in physical adversarial attacks. As \cite{doi:10.24963/ijcai.2021/694} notes, physical adversarial attacks pose a significant threat to the security and reliability of machine learning systems in real-world scenarios. To evaluate the success of such attacks, metrics that consider the attack's success rate, stealthiness, and robustness are essential. 

Building on the concepts of realism and stealth in physical adversarial attacks, one key aspect of evaluating physical adversarial attacks is the measurement of attack success rate. This can be quantified using metrics such as misclassification rate or detection evasion rate, as discussed in \cite{arxiv:1707.08945}. However, these metrics alone do not provide a comprehensive understanding of the attack's effectiveness, as they do not account for the attack's stealthiness or robustness. Stealthiness refers to the ability of the attack to remain undetected by human observers, while robustness refers to the attack's ability to withstand various environmental conditions or defense mechanisms.

To address these limitations, researchers have proposed various evaluation metrics that incorporate stealthiness and robustness. For example, \cite{doi:10.1109/cvpr42600.2020.01426} introduces a metric that evaluates the attack's robustness to different environmental conditions, such as lighting or viewpoint changes. Similarly, \cite{arxiv:2203.03818} proposes a metric that assesses the attack's stealthiness by measuring its ability to remain undetected by human observers. These metrics can be used to evaluate the effectiveness of physical adversarial attacks, which is essential for developing robust defense strategies, as will be discussed in the following section on defense strategies against physical adversarial attacks.

In addition to these metrics, researchers have also explored the use of transferability metrics to evaluate the effectiveness of physical adversarial attacks. Transferability refers to the ability of an attack to be effective against multiple models or systems, as discussed in \cite{doi:10.24963/ijcai.2019/134}. This is particularly important in real-world scenarios, where an attack may need to be effective against multiple systems or models. The development of evaluation metrics that incorporate transferability is crucial for assessing the effectiveness of physical adversarial attacks and developing robust defense strategies.

Despite the progress made in developing evaluation metrics for physical adversarial attacks, there are still several challenges and limitations that need to be addressed. One major challenge is the lack of standardized metrics for evaluating physical adversarial attacks, as highlighted in \cite{arxiv:1902.06705}. This makes it difficult to compare the effectiveness of different attacks or defense mechanisms. Another challenge is the need for more realistic and diverse evaluation scenarios, as discussed in \cite{doi:10.1007/s13042-021-01435-0}. Current evaluation metrics often rely on simplified or simulated scenarios, which may not accurately reflect real-world conditions.

To overcome these challenges, researchers must continue to develop and refine evaluation metrics that can accurately assess the effectiveness and realism of physical adversarial attacks. This may involve incorporating new metrics or techniques, such as those proposed in \cite{arxiv:2305.16494} or \cite{arxiv:2307.15157}. Additionally, there is a need for more comprehensive and standardized evaluation frameworks, as discussed in \cite{doi:10.1109/tpami.2025.3610085}. By addressing these challenges and limitations, researchers can develop more effective evaluation metrics that can help to improve the security and reliability of machine learning systems in real-world scenarios, ultimately informing the development of robust defense strategies against physical adversarial attacks.


\subsection{Defense Strategies Against Physical Adversarial Attacks}


The development of effective defense strategies against physical adversarial attacks is crucial for ensuring the reliability and security of machine learning models in real-world applications. As highlighted in \cite{doi:10.24963/ijcai.2021/694}, physical adversarial attacks can be used to manipulate the inputs to a machine learning model, causing it to make incorrect predictions. To mitigate these attacks, various defense strategies have been proposed, including machine learning-based detection methods and physical countermeasures.

One approach to detecting physical adversarial attacks is to use machine learning-based methods, such as those proposed in \cite{doi:10.1109/tnnls.2021.3105238}. These methods involve training a machine learning model to detect adversarial examples, which can be used to identify and reject physical adversarial attacks. For example, \cite{doi:10.1145/3128572.3140444} demonstrates the use of machine learning-based methods to detect adversarial examples, including those generated using physical adversarial attacks.

In addition to machine learning-based methods, physical countermeasures can also be used to prevent physical adversarial attacks. For example, \cite{arxiv:1807.07769} proposes the use of physical barriers to prevent adversarial objects from being detected by object detectors. Similarly, \cite{arxiv:1907.05418} demonstrates the use of physical countermeasures to prevent adversarial objects from being detected by LiDAR-based autonomous driving systems.

Another approach to defending against physical adversarial attacks is to use adversarial training, which involves training a machine learning model on a dataset that includes adversarial examples. As shown in \cite{arxiv:1803.00404}, adversarial training can be used to improve the robustness of machine learning models to physical adversarial attacks. Additionally, \cite{doi:10.1145/3459637.3482029} demonstrates the use of adversarial training to improve the robustness of deep learning models to physical adversarial attacks.

The use of physical countermeasures and machine learning-based methods can also be combined to provide a more comprehensive defense against physical adversarial attacks. For example, \cite{doi:10.1109/wacv51458.2022.00288} demonstrates the use of physical countermeasures and machine learning-based methods to defend against physical adversarial attacks on semantic segmentation models used in autonomous driving.

In terms of emerging trends and challenges, \cite{doi:10.1007/s11633-019-1211-x} highlights the importance of considering the physical world constraints and limitations when developing defense strategies against physical adversarial attacks. Additionally, \cite{doi:10.1109/iros51168.2021.9636638} demonstrates the need for developing defense strategies that can handle the complexities of multi-modal inputs, such as camera and LiDAR data.

Overall, the development of effective defense strategies against physical adversarial attacks requires a comprehensive approach that considers both machine learning-based methods and physical countermeasures. As highlighted in \cite{doi:10.1145/3702638}, the use of adversarial training, physical countermeasures, and machine learning-based methods can provide a robust defense against physical adversarial attacks. However, further research is needed to address the emerging trends and challenges in this area, including the development of more sophisticated physical adversarial attacks and the need for more robust and generalizable defense strategies.

As noted in \cite{doi:10.1109/cvpr42600.2020.00090}, the use of adversarial examples can actually improve the robustness of machine learning models to physical adversarial attacks. Additionally, \cite{doi:10.1109/iccv.2019.00051} demonstrates the use of adversarial training to improve the robustness of object detection models to physical adversarial attacks.

In conclusion, the development of effective defense strategies against physical adversarial attacks is a critical area of research that requires a comprehensive approach considering both machine learning-based methods and physical countermeasures. By combining these approaches and addressing the emerging trends and challenges in this area, we can develop more robust and reliable machine learning models that can withstand physical adversarial attacks. As highlighted in \cite{doi:10.1007/s11263-023-01934-3}, further research is needed to develop more sophisticated defense strategies that can handle the complexities of physical adversarial attacks.


\subsection{Future Directions and Challenges in Physical Adversarial Attacks}


The field of physical adversarial attacks has witnessed significant advancements in recent years, with various methods being proposed to generate robust and effective attacks against deep learning models in the physical world. As a result, it is essential to address the challenges and open issues that need to be resolved to develop more sophisticated attack methods and robust defense mechanisms. One of the major challenges is the development of more sophisticated attack methods that can effectively evade detection and defense mechanisms \cite{doi:10.1007/s11633-019-1211-x}. This challenge is crucial, as physical adversarial attacks can be used to attack various deep learning models, including object detection models \cite{doi:10.1109/access.2023.3321118}, face recognition models \cite{doi:10.1016/j.patcog.2022.109009}, and semantic segmentation models \cite{doi:10.1016/j.eswa.2025.128967}.

Recent studies have demonstrated the effectiveness of physical adversarial attacks in various scenarios, including digital and physical attacks \cite{doi:10.24963/ijcai.2021/694}. However, the development of more effective attack methods is still an open issue, and further research is needed to address this challenge \cite{doi:10.58496/adsa/2024/011}. To achieve this, researchers must explore new techniques for generating physical adversarial examples, such as using 3D printing and other technologies \cite{doi:10.1109/cvpr.2019.00935}. These examples have been shown to be effective in attacking various deep learning models, including object detection models and face recognition models \cite{doi:10.1016/j.jisa.2020.102694}.

In addition to the development of more sophisticated attack methods, another important direction for future research is the development of robust defense mechanisms against physical adversarial attacks. Various defense mechanisms have been proposed, including detection-based defenses \cite{doi:10.1109/tifs.2023.3278458}, mitigation strategies \cite{arxiv:2002.07405}, and robustness enhancement techniques \cite{arxiv:1712.02779}. However, these defense mechanisms have their own limitations and challenges, and further research is needed to develop more effective and robust defense mechanisms \cite{doi:10.1109/access.2023.3345000}. The development of more effective evaluation metrics and methodologies for physical adversarial attacks is also crucial, as it will enable researchers to assess the effectiveness of defense mechanisms and identify areas for improvement \cite{doi:10.1109/tpami.2024.3430860}.

The development of more effective evaluation metrics and methodologies is essential for advancing the field of physical adversarial attacks. Various evaluation metrics have been proposed, including success rate, detection rate, and robustness metrics \cite{doi:10.1109/tpami.2024.3430860}. However, these metrics have their own limitations and challenges, and further research is needed to develop more effective and robust evaluation metrics and methodologies \cite{arxiv:1902.10899}. By addressing these challenges and developing more sophisticated attack methods, robust defense mechanisms, and effective evaluation metrics and methodologies, researchers can improve the robustness and reliability of deep learning models in the physical world.

In conclusion, the field of physical adversarial attacks is rapidly evolving, with various challenges and open issues that need to be addressed. The development of more sophisticated attack methods, robust defense mechanisms, and effective evaluation metrics and methodologies are some of the key directions for future research \cite{doi:10.1109/tpami.2024.3430860}. As noted in \cite{doi:10.1145/3274895.3274904}, the development of more effective methods for generating physical adversarial examples is crucial for evaluating the robustness of deep learning models in the physical world. Additionally, \cite{doi:10.1109/icip.2019.8803770} highlights the importance of developing robust defense mechanisms against physical adversarial attacks on 3D point cloud classifiers. Overall, the field of physical adversarial attacks requires continued research and development to address the challenges and open issues that exist, and to develop more effective and robust methods for generating and defending against these attacks \cite{doi:10.1007/s10462-024-10841-z}.


\section{Evaluation Metrics and Methodologies}



\subsection{Introduction to Evaluation Metrics}


Evaluating the effectiveness of visual adversarial attacks and defenses in the physical world is crucial for understanding their potential impact on real-world applications. The development of standardized evaluation metrics is essential for comparing different attack and defense strategies and assessing their robustness \cite{doi:10.1145/3128572.3140444}. Existing evaluation metrics for visual adversarial attacks, such as misclassification rates and detection evasion rates, have limitations, including their inability to account for real-world variability and environmental factors \cite{arxiv:1707.08945}. 

To address these limitations, researchers have proposed various evaluation metrics, including robustness metrics, such as adversarial accuracy and certified robustness \cite{arxiv:1902.06705}. These metrics provide a more comprehensive understanding of the robustness of visual adversarial attacks and defenses. For instance, \cite{arxiv:1810.00069} highlights the importance of evaluating the robustness of deep learning models against adversarial attacks. 

Moreover, the evaluation of visual adversarial attacks and defenses in the physical world requires consideration of various factors, including lighting, weather, and viewpoint changes \cite{doi:10.1109/tnnls.2023.3321432}. These factors can significantly impact the effectiveness of visual adversarial attacks and defenses, and their evaluation is crucial for understanding the robustness of these attacks and defenses in real-world scenarios. \cite{doi:10.1109/access.2022.3216291} provides a comprehensive overview of the existing evaluation metrics and methodologies for assessing the robustness of deep learning models against adversarial attacks. 

The development of more robust and generalizable evaluation metrics is an active area of research, with several emerging trends and challenges \cite{doi:10.58496/adsa/2024/011}. For example, \cite{doi:10.1109/tnnls.2025.3561225} highlights the importance of evaluating the robustness of object detection models against adversarial attacks. 

In addition, \cite{doi:10.1007/s10462-024-10841-z} provides a comprehensive overview of the existing methods for generating and countermeasuring adversarial examples in computer vision. 

In conclusion, the evaluation of visual adversarial attacks and defenses in the physical world is a complex task that requires careful consideration of various factors. The development of standardized evaluation metrics, such as robustness metrics, is essential for comparing different attack and defense strategies and assessing their robustness. Emerging trends and challenges in this area include the development of more robust and generalizable evaluation metrics, the evaluation of object detection models against adversarial attacks, and the proposal of novel taxonomic frameworks for categorizing adversarial attacks. 

The importance of evaluating the robustness of deep learning models against adversarial attacks cannot be overstated, and researchers must continue to develop more effective evaluation metrics and methodologies to address the emerging trends and challenges in this area. By doing so, we can better understand the potential impact of visual adversarial attacks and defenses on real-world applications and develop more robust and generalizable defense mechanisms to mitigate these threats. 

Furthermore, \cite{doi:10.1145/3627673.3679750} proposes a novel ViT architecture to achieve a better trade-off between standard accuracy and robustness. 

In the context of facial recognition systems, \cite{doi:10.1109/access.2024.3479949} presents a pilot study introducing an attack-defense framework aimed at enhancing the robustness of facial recognition systems against evasion attacks. 

The evaluation of visual adversarial attacks and defenses in the physical world is an active area of research, with several emerging trends and challenges. The development of more robust and generalizable evaluation metrics and methodologies is crucial for understanding the potential impact of these attacks and defenses on real-world applications. 

In future research, it is essential to consider the development of more robust and generalizable evaluation metrics and methodologies for assessing the robustness of deep learning models against adversarial attacks. This can be achieved by proposing novel taxonomic frameworks for categorizing adversarial attacks, developing more effective evaluation metrics, such as robustness metrics, and evaluating the robustness of object detection models against adversarial attacks. 

Overall, the evaluation of visual adversarial attacks and defenses in the physical world is a complex task that requires careful consideration of various factors. The development of standardized evaluation metrics, such as robustness metrics, is essential for comparing different attack and defense strategies and assessing their robustness. Emerging trends and challenges in this area include the development of more robust and generalizable evaluation metrics, the evaluation of object detection models against adversarial attacks, and the proposal of novel taxonomic frameworks for categorizing adversarial attacks. By continuing to develop more effective evaluation metrics and methodologies, we can create more secure and reliable deep learning models that can withstand the threats posed by visual adversarial attacks and defenses in the physical world.


\subsection{Metrics for Evaluating Attack Success}


Evaluating the success of visual adversarial attacks in the physical world is a complex task that requires careful consideration of various metrics, which is crucial for understanding their potential impact on real-world applications. As discussed in the context of evaluating the effectiveness of visual adversarial attacks and defenses, these metrics can be broadly categorized into three main areas: effectiveness, stealthiness, and robustness. Effectiveness refers to the ability of the attack to achieve its intended goal, such as misclassifying an object or evading detection \cite{arxiv:1712.09665}. Stealthiness, on the other hand, measures how well the attack can avoid being detected by humans or other defense mechanisms \cite{doi:10.1109/cvprw.2019.00012}. Robustness, meanwhile, evaluates the attack's ability to withstand various environmental factors, such as changes in lighting or viewpoint \cite{arxiv:1707.08945}, which is essential for assessing the attack's potential impact in real-world scenarios.

One commonly used metric for evaluating the effectiveness of visual adversarial attacks is the misclassification rate, which measures the proportion of times the attack is successful in misleading the model \cite{doi:10.24963/ijcai.2021/694}. Another metric is the detection evasion rate, which assesses the attack's ability to evade detection by the model \cite{arxiv:1902.06705}. These metrics are essential for understanding the attack's strengths and limitations, and for developing more effective defense mechanisms to counter these attacks. In addition to effectiveness, stealthiness is also an important consideration in evaluating visual adversarial attacks, as it measures the attack's ability to avoid detection by humans or other defense mechanisms. Metrics such as the perceptibility of the adversarial perturbation to human observers can be used to assess the attack's stealthiness \cite{doi:10.1109/tpami.2022.3176760}.

The evaluation of visual adversarial attacks in the physical world also requires consideration of robustness, as it evaluates the attack's ability to withstand various environmental factors. Metrics such as the attack's resistance to changes in lighting or viewpoint can be used to assess robustness \cite{arxiv:2007.16118}. Recent studies have also explored the use of more advanced metrics for evaluating visual adversarial attacks, such as the use of generative models to generate adversarial examples \cite{doi:10.1109/cvpr46437.2021.01167}. These metrics provide a more comprehensive understanding of the attack's strengths and limitations, and can be used to develop more effective defense mechanisms.

In conclusion, evaluating the success of visual adversarial attacks in the physical world requires careful consideration of various metrics, including effectiveness, stealthiness, and robustness. By using a combination of these metrics, researchers can gain a comprehensive understanding of the attack's strengths and limitations, and develop more effective defense mechanisms to counter these attacks. As the field continues to evolve, it is likely that new metrics and evaluation methodologies will be developed to address the increasingly sophisticated nature of visual adversarial attacks \cite{doi:10.1109/tpami.2024.3430860}. Future research directions may include the development of more advanced metrics for evaluating visual adversarial attacks, such as the use of multimodal approaches that combine visual and other modalities, and the development of more effective defense mechanisms, such as the use of adversarial training or robust optimization methods, which will be critical in countering the threat of visual adversarial attacks \cite{doi:10.1007/s11633-019-1211-x}. This will ultimately lead to the development of more secure and reliable deep learning models that can withstand the threats posed by visual adversarial attacks in the physical world, and will be essential for evaluating the effectiveness of defenses against these attacks, as will be discussed in the following section.


\subsection{Metrics for Evaluating Defense Effectiveness}


Evaluating the effectiveness of defenses against visual adversarial attacks in the physical world is a complex task that requires careful consideration of various metrics. At its core, the goal of these defenses is to ensure that machine learning models are robust, accurate, and efficient in the face of adversarial examples. \cite{doi:10.1145/3128572.3140444} have shown that most detection methods can be bypassed, highlighting the need for more robust evaluation metrics. 

One key metric for evaluating defense effectiveness is robustness, which refers to the model's ability to withstand adversarial attacks. \cite{arxiv:1707.07397} proposed a method for synthesizing robust adversarial examples, which can be used to evaluate the robustness of defenses. Adversarial accuracy and certified robustness are also important metrics, as they provide a measure of the model's ability to correctly classify inputs and detect adversarial attacks. \cite{arxiv:1902.06705} discussed the importance of evaluating adversarial robustness and proposed a framework for doing so.

Another important consideration is the trade-off between robustness and accuracy. \cite{arxiv:1511.05644} proposed a method for generating adversarial examples using autoencoders, which can be used to evaluate the robustness of defenses. However, this method may also reduce the model's accuracy, highlighting the need for a balance between robustness and accuracy. \cite{arxiv:1805.06605} proposed a method for defending against adversarial attacks using generative models, which can improve both robustness and accuracy.

Efficiency is also an important metric, as defenses should be able to operate in real-time without significantly impacting the model's performance. \cite{arxiv:2001.05574} proposed a toolbox for generating adversarial examples, which can be used to evaluate the efficiency of defenses. \cite{doi:10.1145/3449726.3463137} proposed a method for generating adversarial examples using evolution strategies, which can be used to evaluate the efficiency of defenses in a black-box setting.

In addition to these metrics, it is also important to consider the physical world constraints and limitations of the defense. \cite{doi:10.1201/9781351251389-8} showed that adversarial examples can be used to attack models in the physical world, highlighting the need for defenses that can operate in this setting. \cite{doi:10.1109/cvpr.2019.00284} proposed a method for improving the transferability of adversarial examples, which can be used to evaluate the effectiveness of defenses in the physical world.

Emerging trends in evaluating defense effectiveness include the use of generative models and evolution strategies to generate adversarial examples. \cite{doi:10.1109/access.2022.3160283} proposed a method for generating adversarial examples using generative models, which can be used to evaluate the robustness of defenses. \cite{doi:10.1186/s41074-019-0053-3} proposed a method for generating adversarial examples using differential evolution, which can be used to evaluate the efficiency of defenses.

In conclusion, evaluating the effectiveness of defenses against visual adversarial attacks in the physical world requires careful consideration of various metrics, including robustness, accuracy, efficiency, and physical world constraints. By using a combination of these metrics and emerging trends, researchers can develop more effective defenses against adversarial attacks. \cite{doi:10.1109/access.2022.3216291} provided a comprehensive survey of methods for improving adversarial robustness, highlighting the need for further research in this area. As the field continues to evolve, it is likely that new metrics and evaluation methods will be developed, providing even more effective defenses against visual adversarial attacks. \cite{doi:10.1109/access.2022.3208131} discussed the importance of evaluating adversarial robustness and proposed a framework for doing so, highlighting the need for further research in this area.


\subsection{Evaluation Methodologies}


Evaluating the effectiveness of visual adversarial attacks and defenses in the physical world is a complex task that requires careful consideration of various factors, including the type of attack, the defense mechanism, and the evaluation methodology. As discussed in the previous section, evaluating the success of visual adversarial attacks in the physical world involves assessing their effectiveness, stealthiness, and robustness. Similarly, when it comes to defenses, the goal is to ensure that machine learning models are robust, accurate, and efficient in the face of adversarial examples. In this context, evaluation methodologies play a crucial role in assessing the robustness of visual adversarial attacks and defenses, as highlighted in \cite{doi:10.24963/ijcai.2021/694}.

Digital testing methodologies, such as simulation-based approaches and dataset-based evaluations, are commonly used to evaluate the effectiveness of visual adversarial attacks and defenses. These methodologies allow for efficient and controlled testing of attacks and defenses, enabling researchers to analyze their performance under various scenarios. For example, \cite{arxiv:1707.08945} uses digital testing to evaluate the robustness of deep learning models against physical-world attacks. However, digital testing has its limitations, as it may not accurately capture the complexities of the physical world, such as environmental factors and variability, which can affect the success of adversarial attacks and the effectiveness of defenses.

Physical testing methodologies, on the other hand, involve testing attacks and defenses in real-world settings, taking into account environmental factors such as lighting, weather, and viewpoint changes. Physical testing provides a more realistic evaluation of the effectiveness of visual adversarial attacks and defenses, as it simulates the conditions under which they will be deployed. \cite{doi:10.1109/cvpr42600.2020.01426} demonstrates the importance of physical testing in evaluating the robustness of autonomous driving systems against physical-world attacks. Furthermore, physical testing can help identify potential vulnerabilities in defenses that may not be apparent in digital testing, as discussed in \cite{doi:10.1109/cvpr42600.2020.00108}.

Hybrid evaluation methodologies combine digital and physical testing approaches to provide a more comprehensive assessment of the effectiveness of visual adversarial attacks and defenses. These methodologies leverage the strengths of both digital and physical testing, enabling researchers to evaluate the performance of attacks and defenses under various scenarios while also considering the complexities of the physical world. \cite{doi:10.1109/eurosp53844.2022.00047} proposes a hybrid evaluation methodology that combines digital and physical testing to evaluate the effectiveness of physical-world attacks on computer vision systems. By using hybrid evaluation methodologies, researchers can develop more robust and generalizable visual adversarial attacks and defenses that can effectively operate in real-world scenarios.

When evaluating the effectiveness of visual adversarial attacks and defenses, it is essential to consider the trade-offs between different evaluation methodologies. Digital testing offers efficiency and control, but may not accurately capture the complexities of the physical world. Physical testing provides a more realistic evaluation, but can be time-consuming and expensive. Hybrid evaluation methodologies offer a balance between these two approaches, enabling researchers to evaluate the effectiveness of visual adversarial attacks and defenses under various scenarios while also considering the complexities of the physical world. As discussed in \cite{doi:10.1109/satml54575.2023.00030}, the evaluation of visual adversarial attacks and defenses is closely tied to the concept of adversarial camouflage, which involves hiding physical-world attacks in natural styles.

In conclusion, the evaluation of visual adversarial attacks and defenses in the physical world requires careful consideration of various factors, including the type of attack, the defense mechanism, and the evaluation methodology. Digital, physical, and hybrid evaluation methodologies each have their strengths and limitations, and the choice of evaluation methodology depends on the specific research question and the trade-offs between efficiency, control, and realism. As the field continues to evolve, it is essential to develop more sophisticated evaluation methodologies that can accurately capture the complexities of the physical world and provide a comprehensive assessment of the effectiveness of visual adversarial attacks and defenses. Future research should focus on developing hybrid evaluation methodologies that combine the strengths of digital and physical testing, and on evaluating the effectiveness of visual adversarial attacks and defenses in real-world settings, taking into account factors such as environmental variability and adversarial camouflage, as discussed in \cite{arxiv:2203.03818} and \cite{doi:10.1109/access.2024.3446834}. This will be further discussed in the following section on evaluation methodologies for visual adversarial attacks and defenses in the physical world.


\subsection{Real-World Considerations and Constraints}


The evaluation of visual adversarial attacks and defenses in the physical world is a complex task that requires consideration of various real-world factors and constraints. As discussed in \cite{doi:10.24963/ijcai.2021/694}, the physical world poses unique challenges to the development and evaluation of visual adversarial attacks and defenses, including environmental factors, variability, and scalability. In this subsection, we delve into these considerations and constraints, providing a comprehensive analysis of their impact on the evaluation of visual adversarial attacks and defenses.

Environmental factors, such as lighting, weather, and viewpoint changes, can significantly affect the success of visual adversarial attacks and the effectiveness of defenses \cite{arxiv:1807.07769}. For instance, a visual adversarial attack that is effective in a laboratory setting may not be effective in a real-world setting with varying lighting conditions. Similarly, defenses that are designed to detect adversarial examples in a controlled environment may not be effective in a real-world setting with diverse environmental factors. As noted in \cite{arxiv:1707.08945}, the development of robust visual adversarial attacks and defenses requires careful consideration of these environmental factors.

Variability is another critical consideration in the evaluation of visual adversarial attacks and defenses. Real-world scenarios often involve variability in object appearance, pose, and context, which can affect the success of visual adversarial attacks and the effectiveness of defenses \cite{doi:10.3390/cancers15174228}. For example, a visual adversarial attack that is designed to target a specific object may not be effective if the object is viewed from a different angle or in a different context. Defenses that are designed to detect adversarial examples in a specific context may not be effective in other contexts. As discussed in \cite{doi:10.1016/j.eng.2019.12.012}, the development of robust visual adversarial attacks and defenses requires careful consideration of variability in real-world scenarios.

Scalability is also an essential consideration in the evaluation of visual adversarial attacks and defenses. Real-world applications often involve large-scale deployments, which can affect the success of visual adversarial attacks and the effectiveness of defenses \cite{doi:10.1109/access.2022.3208131}. For instance, a visual adversarial attack that is effective in a small-scale setting may not be effective in a large-scale setting with multiple objects and complex scenes. Defenses that are designed to detect adversarial examples in a small-scale setting may not be effective in a large-scale setting. As noted in \cite{doi:10.1007/s11263-023-01934-3}, the development of robust visual adversarial attacks and defenses requires careful consideration of scalability in real-world applications.

In addition to these considerations, the evaluation of visual adversarial attacks and defenses in the physical world also requires careful consideration of the trade-offs between different evaluation metrics and methodologies. As discussed in \cite{doi:10.3390/bdcc8010008}, different evaluation metrics and methodologies can provide different insights into the effectiveness of visual adversarial attacks and defenses. For example, metrics that focus on the success rate of visual adversarial attacks may not provide a complete picture of their effectiveness in real-world scenarios. Similarly, methodologies that focus on detecting adversarial examples in a specific context may not be effective in other contexts. As noted in \cite{doi:10.1145/3702638}, the development of robust visual adversarial attacks and defenses requires careful consideration of the trade-offs between different evaluation metrics and methodologies.

In conclusion, the evaluation of visual adversarial attacks and defenses in the physical world is a complex task that requires careful consideration of various real-world factors and constraints. Environmental factors, variability, and scalability are critical considerations that can affect the success of visual adversarial attacks and the effectiveness of defenses. The development of robust visual adversarial attacks and defenses requires careful consideration of these factors, as well as the trade-offs between different evaluation metrics and methodologies. As noted in \cite{doi:10.1016/j.sysarc.2020.101766}, future research should focus on developing more robust and generalizable visual adversarial attacks and defenses that can effectively operate in real-world scenarios. By providing a comprehensive analysis of the real-world considerations and constraints that affect the evaluation of visual adversarial attacks and defenses, this subsection aims to contribute to the development of more robust and reliable visual adversarial attacks and defenses for real-world applications. Further research in this area is necessary to fully understand the implications of visual adversarial attacks and defenses in the physical world, as discussed in \cite{arxiv:2301.10766}.


\subsection{Future Directions and Open Challenges}


The development of evaluation metrics and methodologies for visual adversarial attacks and defenses in the physical world is a rapidly evolving field, with significant implications for the security and reliability of machine learning models in real-world applications \cite{doi:10.1109/tpami.2024.3430860}. As discussed in the previous subsection, the evaluation of visual adversarial attacks and defenses in the physical world is a complex task that requires careful consideration of various real-world factors and constraints, including environmental factors, variability, and scalability. Building on this discussion, several future directions and open challenges have emerged in the development of evaluation metrics and methodologies for visual adversarial attacks and defenses in the physical world \cite{doi:10.24963/ijcai.2021/694}. 

One of the primary challenges is the development of more robust and generalizable evaluation metrics that can account for the variability and unpredictability of real-world environments \cite{arxiv:1707.08945}. Current evaluation metrics, such as misclassification rates and detection evasion rates, are often limited in their ability to capture the full range of potential attack scenarios and defense strategies \cite{doi:10.1109/access.2018.2807385}. To address this challenge, researchers must develop evaluation metrics that can effectively capture the complexities of real-world environments and provide a comprehensive understanding of the effectiveness of visual adversarial attacks and defenses.

Another significant challenge is the need for more efficient and scalable evaluation methodologies that can accommodate large-scale, real-world deployments \cite{doi:10.1109/icpr48806.2021.9412236}. Digital testing methodologies, such as simulation-based approaches and dataset-based evaluations, have been widely used, but they often fail to capture the complexities of real-world environments \cite{doi:10.1109/access.2023.3321118}. Physical testing methodologies, such as laboratory-based experiments and real-world field tests, are more effective but often time-consuming and resource-intensive \cite{doi:10.1007/978-3-030-58589-1_3}. Hybrid evaluation methodologies that combine digital and physical testing approaches may offer a more comprehensive and efficient solution \cite{doi:10.1109/tgrs.2022.3225306}.

The development of more sophisticated attack methods and robust defense mechanisms is also an open challenge \cite{doi:10.1007/s11633-019-1211-x}. Recent studies have shown that adversarial patches can be designed to evade detection by object detectors \cite{doi:10.1109/cvprw.2019.00012}, while other studies have demonstrated the effectiveness of defense mechanisms, such as adversarial training and defensive distillation \cite{doi:10.1016/j.eng.2019.12.012}. However, the cat-and-mouse game between attackers and defenders is ongoing, and more research is needed to develop effective and robust defense strategies \cite{arxiv:2007.16118}.

The application of visual adversarial attacks and defenses in various domains, such as security, healthcare, and transportation, is another area of ongoing research \cite{doi:10.1109/jstars.2024.3422377}. The use of deep learning models in these domains has significant potential benefits, but it also raises concerns about the potential risks and vulnerabilities of these models to adversarial attacks \cite{doi:10.1109/cvprw56347.2022.00034}. More research is needed to develop effective defense strategies and evaluation methodologies that can address the unique challenges and requirements of these domains \cite{doi:10.1007/s10462-024-10841-z}.

In conclusion, the development of evaluation metrics and methodologies for visual adversarial attacks and defenses in the physical world is a complex and ongoing challenge \cite{doi:10.58496/adsa/2024/011}. As researchers continue to explore the complexities of visual adversarial attacks, it is essential to develop more robust and generalizable evaluation metrics, efficient and scalable evaluation methodologies, and sophisticated attack methods and robust defense mechanisms \cite{doi:10.1109/vrw62533.2024.00284}. The application of visual adversarial attacks and defenses in various domains also requires more research and development of effective defense strategies and evaluation methodologies \cite{arxiv:2409.20426}. By addressing these challenges and open questions, researchers can help to ensure the security and reliability of machine learning models in real-world applications, which will be further discussed in the following subsection on the future directions and open challenges in the development of visual adversarial attacks and defenses in the physical world \cite{doi:10.1109/tpami.2024.3430860}.


\section{Conclusion}


The field of visual adversarial attacks and defenses in the physical world has experienced significant growth in recent years, with a multitude of research efforts focused on understanding and mitigating the vulnerabilities of deep learning models in real-world scenarios \cite{doi:10.1145/3128572.3140444}. A key aspect of this research is the development of effective attack methods, such as the creation of physical adversarial patches that can deceive object detectors in various environments \cite{arxiv:1707.08945}. These patches can be designed to be robust against different lighting conditions, viewpoints, and other environmental factors, making them a formidable challenge for defense mechanisms \cite{arxiv:1712.09665}.

One of the primary challenges in defending against visual adversarial attacks is the need for robust and generalizable defense mechanisms \cite{arxiv:1902.06705}. Current defense strategies, such as adversarial training and input preprocessing, have shown promise in improving the robustness of deep learning models, but they are often limited in their effectiveness against strong attacks \cite{arxiv:1810.00069}. Furthermore, the development of certified defenses, which can provide formal guarantees of robustness, is an active area of research \cite{arxiv:2003.06693}.

In addition to the development of defense mechanisms, there is a growing need for standardized evaluation metrics and methodologies for assessing the effectiveness of visual adversarial attacks and defenses \cite{arxiv:2010.09670}. This includes the development of robustness metrics, such as adversarial accuracy and certified robustness, which can provide a comprehensive understanding of a model's vulnerability to attacks. Moreover, the creation of benchmark datasets, such as APRICOT, can facilitate the evaluation and comparison of different attack and defense methods \cite{doi:10.1007/978-3-030-58589-1_3}.

The physical world imposes unique challenges on visual adversarial attacks and defenses, such as the need for realism and stealthiness \cite{doi:10.1109/tnnls.2023.3321432}. Attackers must design patches that are not only effective but also imperceptible to humans, making them more difficult to detect \cite{doi:10.1109/tpami.2022.3176760}. Furthermore, the development of attacks that can be implemented in the physical world, such as the use of projectors or 3D printing, requires careful consideration of the physical properties of the attack \cite{doi:10.1016/j.eswa.2025.128967}.

Future research directions in the field of visual adversarial attacks and defenses include the development of more robust and generalizable defense mechanisms, the creation of standardized evaluation metrics and methodologies, and the exploration of new attack methods and scenarios \cite{doi:10.1109/access.2022.3216291}. Additionally, the development of attacks and defenses for specific applications, such as autonomous vehicles and surveillance systems, is an important area of research \cite{doi:10.1109/access.2023.3321118}. Overall, the field of visual adversarial attacks and defenses in the physical world is rapidly evolving, with new challenges and opportunities emerging continuously \cite{doi:10.58496/adsa/2024/011}. As such, it is essential to continue researching and developing effective defense mechanisms to ensure the security and reliability of deep learning models in real-world scenarios \cite{doi:10.59247/csol.v2i1.73}.

\begin{thebibliography}{129}
\setlength{\itemsep}{2pt}

\bibitem{doi:10.24963/ijcai.2021/694}
\newblock \emph{Adversarial Examples in Physical World}.
\newblock DOI:10.24963/ijcai.2021/694, 2021.
\newblock \url{https://doi.org/10.24963/ijcai.2021/694}

\bibitem{doi:10.1007/s11633-019-1211-x}
\newblock \emph{Adversarial Attacks and Defenses in Images, Graphs and Text: A Review}.
\newblock DOI:10.1007/s11633-019-1211-x, 2020.
\newblock \url{https://doi.org/10.1007/s11633-019-1211-x}

\bibitem{doi:10.1145/3128572.3140444}
\newblock \emph{Adversarial Examples Are Not Easily Detected: Bypassing Ten Detection Methods}.
\newblock DOI:10.1145/3128572.3140444, 2017.
\newblock \url{https://doi.org/10.1145/3128572.3140444}

\bibitem{arxiv:1707.08945}
\newblock \emph{Robust Physical-World Attacks on Deep Learning Models}.
\newblock arXiv:1707.08945, 2017.
\newblock \url{https://arxiv.org/abs/1707.08945}

\bibitem{arxiv:1712.09665}
\newblock \emph{Adversarial Patch}.
\newblock arXiv:1712.09665, 2017.
\newblock \url{https://arxiv.org/abs/1712.09665}

\bibitem{doi:10.1007/978-3-030-10925-7_4}
\newblock \emph{ShapeShifter: Robust Physical Adversarial Attack on Faster R-CNN Object Detector}.
\newblock DOI:10.1007/978-3-030-10925-7\_4, 2019.
\newblock \url{https://doi.org/10.1007/978-3-030-10925-7_4}

\bibitem{arxiv:1906.11897}
\newblock \emph{On Physical Adversarial Patches for Object Detection}.
\newblock arXiv:1906.11897, 2019.
\newblock \url{https://arxiv.org/abs/1906.11897}

\bibitem{doi:10.1007/978-3-030-58574-7_41}
\newblock \emph{PatchAttack: A Black-box Texture-based Attack with Reinforcement Learning}.
\newblock DOI:10.1007/978-3-030-58574-7\_41, 2020.
\newblock \url{https://doi.org/10.1007/978-3-030-58574-7_41}

\bibitem{doi:10.1109/tnnls.2023.3321432}
\newblock \emph{Physical Adversarial Attacks for Surveillance: A Survey}.
\newblock DOI:10.1109/tnnls.2023.3321432, 2023.
\newblock \url{https://doi.org/10.1109/tnnls.2023.3321432}

\bibitem{doi:10.1109/tpami.2024.3430860}
\newblock \emph{Physical Adversarial Attack Meets Computer Vision: A Decade Survey}.
\newblock DOI:10.1109/tpami.2024.3430860, 2024.
\newblock \url{https://doi.org/10.1109/tpami.2024.3430860}

\bibitem{doi:10.1109/iccv48922.2021.01542}
\newblock \emph{Robustness and Generalization via Generative Adversarial Training}.
\newblock DOI:10.1109/iccv48922.2021.01542, 2021.
\newblock \url{https://doi.org/10.1109/iccv48922.2021.01542}

\bibitem{arxiv:2003.06693}
\newblock \emph{Certified Defenses for Adversarial Patches}.
\newblock arXiv:2003.06693, 2020.
\newblock \url{https://arxiv.org/abs/2003.06693}

\bibitem{arxiv:1902.06705}
\newblock \emph{On Evaluating Adversarial Robustness}.
\newblock arXiv:1902.06705, 2019.
\newblock \url{https://arxiv.org/abs/1902.06705}

\bibitem{doi:10.1109/access.2022.3216291}
\newblock \emph{A Survey on Efficient Methods for Adversarial Robustness}.
\newblock DOI:10.1109/access.2022.3216291, 2022.
\newblock \url{https://doi.org/10.1109/access.2022.3216291}

\bibitem{doi:10.1016/j.ins.2023.03.084}
\newblock \emph{3D Adversarial Attacks Beyond Point Cloud}.
\newblock DOI:10.1016/j.ins.2023.03.084, 2023.
\newblock \url{https://doi.org/10.1016/j.ins.2023.03.084}

\bibitem{doi:10.1016/j.eng.2019.12.012}
\newblock \emph{Adversarial Attacks and Defenses in Deep Learning}.
\newblock DOI:10.1016/j.eng.2019.12.012, 2020.
\newblock \url{https://doi.org/10.1016/j.eng.2019.12.012}

\bibitem{arxiv:1810.00069}
\newblock \emph{Adversarial Attacks and Defences: A Survey}.
\newblock arXiv:1810.00069, 2018.
\newblock \url{https://arxiv.org/abs/1810.00069}

\bibitem{arxiv:2002.07405}
\newblock \emph{Deflecting Adversarial Attacks}.
\newblock arXiv:2002.07405, 2020.
\newblock \url{https://arxiv.org/abs/2002.07405}

\bibitem{arxiv:1807.07769}
\newblock \emph{Physical Adversarial Examples for Object Detectors}.
\newblock arXiv:1807.07769, 2018.
\newblock \url{https://arxiv.org/abs/1807.07769}

\bibitem{doi:10.1109/access.2018.2807385}
\newblock \emph{Threat of Adversarial Attacks on Deep Learning in Computer Vision: A Survey}.
\newblock DOI:10.1109/access.2018.2807385, 2018.
\newblock \url{https://doi.org/10.1109/access.2018.2807385}

\bibitem{doi:10.1049/cit2.12028}
\newblock \emph{A survey on adversarial attacks and defences}.
\newblock DOI:10.1049/cit2.12028, 2021.
\newblock \url{https://doi.org/10.1049/cit2.12028}

\bibitem{doi:10.1016/j.cose.2020.102092}
\newblock \emph{A Survey of Machine Learning Techniques in Adversarial Image Forensics}.
\newblock DOI:10.1016/j.cose.2020.102092, 2020.
\newblock \url{https://doi.org/10.1016/j.cose.2020.102092}

\bibitem{doi:10.58496/adsa/2024/011}
\newblock \emph{Adversarial Attacks in Machine Learning: Key Insights and Defense Approaches}.
\newblock DOI:10.58496/adsa/2024/011, 2024.
\newblock \url{https://doi.org/10.58496/adsa/2024/011}

\bibitem{doi:10.1109/lra.2022.3189783}
\newblock \emph{Physical Adversarial Attack on a Robotic Arm}.
\newblock DOI:10.1109/lra.2022.3189783, 2022.
\newblock \url{https://doi.org/10.1109/lra.2022.3189783}

\bibitem{doi:10.1109/tnnls.2025.3561225}
\newblock \emph{A Survey and Evaluation of Adversarial Attacks in Object Detection}.
\newblock DOI:10.1109/tnnls.2025.3561225, 2025.
\newblock \url{https://doi.org/10.1109/tnnls.2025.3561225}

\bibitem{arxiv:1806.02299}
\newblock \emph{DPatch: An Adversarial Patch Attack on Object Detectors}.
\newblock arXiv:1806.02299, 2018.
\newblock \url{https://arxiv.org/abs/1806.02299}

\bibitem{arxiv:1707.07397}
\newblock \emph{Synthesizing Robust Adversarial Examples}.
\newblock arXiv:1707.07397, 2017.
\newblock \url{https://arxiv.org/abs/1707.07397}

\bibitem{doi:10.1109/cvpr.2019.00284}
\newblock \emph{Improving Transferability of Adversarial Examples With Input Diversity}.
\newblock DOI:10.1109/cvpr.2019.00284, 2019.
\newblock \url{https://doi.org/10.1109/cvpr.2019.00284}

\bibitem{doi:10.1109/tip.2021.3092822}
\newblock \emph{Toward Visual Distortion in Black-Box Attacks}.
\newblock DOI:10.1109/tip.2021.3092822, 2021.
\newblock \url{https://doi.org/10.1109/tip.2021.3092822}

\bibitem{arxiv:1805.06605}
\newblock \emph{Defense-GAN: Protecting Classifiers Against Adversarial Attacks Using Generative Models}.
\newblock arXiv:1805.06605, 2018.
\newblock \url{https://arxiv.org/abs/1805.06605}

\bibitem{doi:10.1007/s13042-021-01435-0}
\newblock \emph{Really natural adversarial examples}.
\newblock DOI:10.1007/s13042-021-01435-0, 2021.
\newblock \url{https://doi.org/10.1007/s13042-021-01435-0}

\bibitem{doi:10.1609/aaai.v32i1.11828}
\newblock \emph{Detecting Adversarial Examples Through Image Transformation}.
\newblock DOI:10.1609/aaai.v32i1.11828, 2018.
\newblock \url{https://doi.org/10.1609/aaai.v32i1.11828}

\bibitem{doi:10.24963/ijcai.2019/134}
\newblock \emph{Transferable Adversarial Attacks for Image and Video Object Detection}.
\newblock DOI:10.24963/ijcai.2019/134, 2019.
\newblock \url{https://doi.org/10.24963/ijcai.2019/134}

\bibitem{doi:10.1186/s41074-019-0053-3}
\newblock \emph{Attacking convolutional neural network using differential evolution}.
\newblock DOI:10.1186/s41074-019-0053-3, 2019.
\newblock \url{https://doi.org/10.1186/s41074-019-0053-3}

\bibitem{doi:10.1609/aaai.v34i07.6918}
\newblock \emph{Heuristic Black-box Adversarial Attacks on Video Recognition Models}.
\newblock DOI:10.1609/aaai.v34i07.6918, 2020.
\newblock \url{https://doi.org/10.1609/aaai.v34i07.6918}

\bibitem{doi:10.1007/s00521-021-06812-y}
\newblock \emph{Understanding deep learning defenses against adversarial examples through visualizations for dynamic risk assessment}.
\newblock DOI:10.1007/s00521-021-06812-y, 2022.
\newblock \url{https://doi.org/10.1007/s00521-021-06812-y}

\bibitem{arxiv:1511.07528}
\newblock \emph{The Limitations of Deep Learning in Adversarial Settings}.
\newblock arXiv:1511.07528, 2015.
\newblock \url{https://arxiv.org/abs/1511.07528}

\bibitem{doi:10.1201/9781351251389-8}
\newblock \emph{Adversarial examples in the physical world}.
\newblock DOI:10.1201/9781351251389-8, 2018.
\newblock \url{https://doi.org/10.1201/9781351251389-8}

\bibitem{arxiv:1810.05162}
\newblock \emph{Characterizing Adversarial Examples Based on Spatial Consistency Information for Semantic Segmentation}.
\newblock arXiv:1810.05162, 2018.
\newblock \url{https://arxiv.org/abs/1810.05162}

\bibitem{doi:10.1145/3133956.3134057}
\newblock \emph{MagNet: A Two-Pronged Defense against Adversarial Examples}.
\newblock DOI:10.1145/3133956.3134057, 2017.
\newblock \url{https://doi.org/10.1145/3133956.3134057}

\bibitem{doi:10.1109/iros51168.2021.9636638}
\newblock \emph{Adversarial Attacks on Camera-LiDAR Models for 3D Car Detection}.
\newblock DOI:10.1109/iros51168.2021.9636638, 2021.
\newblock \url{https://doi.org/10.1109/iros51168.2021.9636638}

\bibitem{arxiv:1904.06347}
\newblock \emph{Unrestricted Adversarial Examples via Semantic Manipulation}.
\newblock arXiv:1904.06347, 2019.
\newblock \url{https://arxiv.org/abs/1904.06347}

\bibitem{doi:10.3390/bdcc8010008}
\newblock \emph{Evaluating the Robustness of Deep Learning Models against Adversarial Attacks: An Analysis with FGSM, PGD and CW}.
\newblock DOI:10.3390/bdcc8010008, 2024.
\newblock \url{https://doi.org/10.3390/bdcc8010008}

\bibitem{doi:10.1109/wacv51458.2022.00144}
\newblock \emph{Adversarial Robustness of Deep Sensor Fusion Models}.
\newblock DOI:10.1109/wacv51458.2022.00144, 2022.
\newblock \url{https://doi.org/10.1109/wacv51458.2022.00144}

\bibitem{doi:10.1007/978-3-030-58589-1_3}
\newblock \emph{APRICOT: A Dataset of Physical Adversarial Attacks on Object Detection}.
\newblock DOI:10.1007/978-3-030-58589-1\_3, 2020.
\newblock \url{https://doi.org/10.1007/978-3-030-58589-1_3}

\bibitem{doi:10.1109/access.2022.3208131}
\newblock \emph{Adversarial Deep Learning: A Survey on Adversarial Attacks and Defense Mechanisms on Image Classification}.
\newblock DOI:10.1109/access.2022.3208131, 2022.
\newblock \url{https://doi.org/10.1109/access.2022.3208131}

\bibitem{arxiv:2010.09670}
\newblock \emph{RobustBench: a standardized adversarial robustness benchmark}.
\newblock arXiv:2010.09670, 2020.
\newblock \url{https://arxiv.org/abs/2010.09670}

\bibitem{doi:10.1109/cvprw.2019.00012}
\newblock \emph{Fooling Automated Surveillance Cameras: Adversarial Patches to Attack Person Detection}.
\newblock DOI:10.1109/cvprw.2019.00012, 2019.
\newblock \url{https://doi.org/10.1109/cvprw.2019.00012}

\bibitem{doi:10.1109/tcsvt.2023.3341170}
\newblock \emph{Adversarial Attacks on Video Object Segmentation With Hard Region Discovery}.
\newblock DOI:10.1109/tcsvt.2023.3341170, 2023.
\newblock \url{https://doi.org/10.1109/tcsvt.2023.3341170}

\bibitem{doi:10.1007/s44267-024-00066-7}
\newblock \emph{Effectiveness assessment of recent large vision-language models}.
\newblock DOI:10.1007/s44267-024-00066-7, 2024.
\newblock \url{https://doi.org/10.1007/s44267-024-00066-7}

\bibitem{doi:10.1109/cvpr46437.2021.01167}
\newblock \emph{Improving Transferability of Adversarial Patches on Face Recognition with Generative Models}.
\newblock DOI:10.1109/cvpr46437.2021.01167, 2021.
\newblock \url{https://doi.org/10.1109/cvpr46437.2021.01167}

\bibitem{doi:10.3390/s23020853}
\newblock \emph{Adversarial Patch Attacks on Deep-Learning-Based Face Recognition Systems Using Generative Adversarial Networks}.
\newblock DOI:10.3390/s23020853, 2023.
\newblock \url{https://doi.org/10.3390/s23020853}

\bibitem{arxiv:1606.03498}
\newblock \emph{Improved Techniques for Training GANs}.
\newblock arXiv:1606.03498, 2016.
\newblock \url{https://arxiv.org/abs/1606.03498}

\bibitem{arxiv:1511.05644}
\newblock \emph{Adversarial Autoencoders}.
\newblock arXiv:1511.05644, 2015.
\newblock \url{https://arxiv.org/abs/1511.05644}

\bibitem{arxiv:1901.10513}
\newblock \emph{Adversarial Examples Are a Natural Consequence of Test Error in Noise}.
\newblock arXiv:1901.10513, 2019.
\newblock \url{https://arxiv.org/abs/1901.10513}

\bibitem{arxiv:2001.05574}
\newblock \emph{Advbox: a toolbox to generate adversarial examples that fool neural networks}.
\newblock arXiv:2001.05574, 2020.
\newblock \url{https://arxiv.org/abs/2001.05574}

\bibitem{doi:10.1109/eurosp53844.2022.00047}
\newblock \emph{GRAPHITE: Generating Automatic Physical Examples for Machine-Learning Attacks on Computer Vision Systems}.
\newblock DOI:10.1109/eurosp53844.2022.00047, 2022.
\newblock \url{https://doi.org/10.1109/eurosp53844.2022.00047}

\bibitem{doi:10.1109/cvpr42600.2020.00090}
\newblock \emph{Adversarial Examples Improve Image Recognition}.
\newblock DOI:10.1109/cvpr42600.2020.00090, 2020.
\newblock \url{https://doi.org/10.1109/cvpr42600.2020.00090}

\bibitem{doi:10.1145/3398394}
\newblock \emph{Adversarial Examples on Object Recognition}.
\newblock DOI:10.1145/3398394, 2020.
\newblock \url{https://doi.org/10.1145/3398394}

\bibitem{doi:10.3390/electronics13132458}
\newblock \emph{Research and Application of the Median Filtering Method in Enhancing the Imperceptibility of Perturbations in Adversarial Examples}.
\newblock DOI:10.3390/electronics13132458, 2024.
\newblock \url{https://doi.org/10.3390/electronics13132458}

\bibitem{doi:10.1109/cvpr42600.2020.01426}
\newblock \emph{PhysGAN: Generating Physical-World-Resilient Adversarial Examples for Autonomous Driving}.
\newblock DOI:10.1109/cvpr42600.2020.01426, 2020.
\newblock \url{https://doi.org/10.1109/cvpr42600.2020.01426}

\bibitem{doi:10.1109/cvpr42600.2020.00108}
\newblock \emph{Adversarial Camouflage: Hiding Physical-World Attacks With Natural Styles}.
\newblock DOI:10.1109/cvpr42600.2020.00108, 2020.
\newblock \url{https://doi.org/10.1109/cvpr42600.2020.00108}

\bibitem{arxiv:2203.03818}
\newblock \emph{Shadows can be Dangerous: Stealthy and Effective Physical-world Adversarial Attack by Natural Phenomenon}.
\newblock arXiv:2203.03818, 2022.
\newblock \url{https://arxiv.org/abs/2203.03818}

\bibitem{doi:10.1109/tnnls.2021.3105238}
\newblock \emph{Detecting Adversarial Examples by Input Transformations, Defense Perturbations, and Voting}.
\newblock DOI:10.1109/tnnls.2021.3105238, 2021.
\newblock \url{https://doi.org/10.1109/tnnls.2021.3105238}

\bibitem{arxiv:1904.13000}
\newblock \emph{Adversarial Training and Robustness for Multiple Perturbations}.
\newblock arXiv:1904.13000, 2019.
\newblock \url{https://arxiv.org/abs/1904.13000}

\bibitem{doi:10.1007/s10462-024-10841-z}
\newblock \emph{Generation and Countermeasures of adversarial examples on vision: a survey}.
\newblock DOI:10.1007/s10462-024-10841-z, 2024.
\newblock \url{https://doi.org/10.1007/s10462-024-10841-z}

\bibitem{doi:10.1109/wacv57701.2024.00393}
\newblock \emph{Exploring Adversarial Robustness of Vision Transformers in the Spectral Perspective}.
\newblock DOI:10.1109/wacv57701.2024.00393, 2024.
\newblock \url{https://doi.org/10.1109/wacv57701.2024.00393}

\bibitem{arxiv:2307.15157}
\newblock \emph{R-LPIPS: An Adversarially Robust Perceptual Similarity Metric}.
\newblock arXiv:2307.15157, 2023.
\newblock \url{https://arxiv.org/abs/2307.15157}

\bibitem{arxiv:2305.16494}
\newblock \emph{Diffusion-Based Adversarial Sample Generation for Improved Stealthiness and Controllability}.
\newblock arXiv:2305.16494, 2023.
\newblock \url{https://arxiv.org/abs/2305.16494}

\bibitem{doi:10.3390/electronics10172132}
\newblock \emph{A Survey on Adversarial Deep Learning Robustness in Medical Image Analysis}.
\newblock DOI:10.3390/electronics10172132, 2021.
\newblock \url{https://doi.org/10.3390/electronics10172132}

\bibitem{doi:10.1145/3702638}
\newblock \emph{Survey on Adversarial Attack and Defense for Medical Image Analysis: Methods and Challenges}.
\newblock DOI:10.1145/3702638, 2024.
\newblock \url{https://doi.org/10.1145/3702638}

\bibitem{doi:10.1016/j.sysarc.2020.101766}
\newblock \emph{Attacking Vision-based Perception in End-to-End Autonomous Driving Models}.
\newblock DOI:10.1016/j.sysarc.2020.101766, 2020.
\newblock \url{https://doi.org/10.1016/j.sysarc.2020.101766}

\bibitem{arxiv:1803.00404}
\newblock \emph{Deep Defense: Training DNNs with Improved Adversarial Robustness}.
\newblock arXiv:1803.00404, 2018.
\newblock \url{https://arxiv.org/abs/1803.00404}

\bibitem{doi:10.1007/s11263-023-01934-3}
\newblock \emph{A Comprehensive Study of the Robustness for LiDAR-Based 3D Object Detectors Against Adversarial Attacks}.
\newblock DOI:10.1007/s11263-023-01934-3, 2023.
\newblock \url{https://doi.org/10.1007/s11263-023-01934-3}

\bibitem{doi:10.1109/access.2023.3345000}
\newblock \emph{Adversarial Attacks and Defenses on 3D Point Cloud Classification: A Survey}.
\newblock DOI:10.1109/access.2023.3345000, 2023.
\newblock \url{https://doi.org/10.1109/access.2023.3345000}

\bibitem{arxiv:2210.07540}
\newblock \emph{When Adversarial Training Meets Vision Transformers: Recipes from Training to Architecture}.
\newblock arXiv:2210.07540, 2022.
\newblock \url{https://arxiv.org/abs/2210.07540}

\bibitem{doi:10.1109/cvprw56347.2022.00034}
\newblock \emph{Adversarial Machine Learning Attacks Against Video Anomaly Detection Systems}.
\newblock DOI:10.1109/cvprw56347.2022.00034, 2022.
\newblock \url{https://doi.org/10.1109/cvprw56347.2022.00034}

\bibitem{doi:10.1109/cvpr.2019.00935}
\newblock \emph{Generating 3D Adversarial Point Clouds}.
\newblock DOI:10.1109/cvpr.2019.00935, 2019.
\newblock \url{https://doi.org/10.1109/cvpr.2019.00935}

\bibitem{arxiv:1907.05418}
\newblock \emph{Adversarial Objects Against LiDAR-Based Autonomous Driving Systems}.
\newblock arXiv:1907.05418, 2019.
\newblock \url{https://arxiv.org/abs/1907.05418}

\bibitem{doi:10.1109/access.2023.3321118}
\newblock \emph{Physical Adversarial Attacks for Camera-Based Smart Systems: Current Trends, Categorization, Applications, Research Challenges, and Future Outlook}.
\newblock DOI:10.1109/access.2023.3321118, 2023.
\newblock \url{https://doi.org/10.1109/access.2023.3321118}

\bibitem{doi:10.1145/3274895.3274904}
\newblock \emph{Adversarial examples in remote sensing}.
\newblock DOI:10.1145/3274895.3274904, 2018.
\newblock \url{https://doi.org/10.1145/3274895.3274904}

\bibitem{arxiv:1902.10899}
\newblock \emph{Adversarial Attack and Defense on Point Sets}.
\newblock arXiv:1902.10899, 2019.
\newblock \url{https://arxiv.org/abs/1902.10899}

\bibitem{doi:10.1016/j.eswa.2025.128967}
\newblock \emph{AdvReal: Physical adversarial patch generation framework for security evaluation of object detection systems}.
\newblock DOI:10.1016/j.eswa.2025.128967, 2025.
\newblock \url{https://doi.org/10.1016/j.eswa.2025.128967}

\bibitem{arxiv:2409.20426}
\newblock \emph{Navigating Threats: A Survey of Physical Adversarial Attacks on LiDAR Perception Systems in Autonomous Vehicles}.
\newblock arXiv:2409.20426, 2024.
\newblock \url{https://arxiv.org/abs/2409.20426}

\bibitem{arxiv:2005.10884}
\newblock \emph{PatchGuard: A Provably Robust Defense against Adversarial Patches via Small Receptive Fields and Masking}.
\newblock arXiv:2005.10884, 2020.
\newblock \url{https://arxiv.org/abs/2005.10884}

\bibitem{doi:10.1109/cvpr52688.2022.01455}
\newblock \emph{Segment and Complete: Defending Object Detectors against Adversarial Patch Attacks with Robust Patch Detection}.
\newblock DOI:10.1109/cvpr52688.2022.01455, 2022.
\newblock \url{https://doi.org/10.1109/cvpr52688.2022.01455}

\bibitem{doi:10.1109/access.2024.3435347}
\newblock \emph{Adversarial Robustness of Vision Transformers Versus Convolutional Neural Networks}.
\newblock DOI:10.1109/access.2024.3435347, 2024.
\newblock \url{https://doi.org/10.1109/access.2024.3435347}

\bibitem{doi:10.1109/cvpr52688.2022.01480}
\newblock \emph{Give Me Your Attention: Dot-Product Attention Considered Harmful for Adversarial Patch Robustness}.
\newblock DOI:10.1109/cvpr52688.2022.01480, 2022.
\newblock \url{https://doi.org/10.1109/cvpr52688.2022.01480}

\bibitem{doi:10.1109/cvpr52729.2023.01187}
\newblock \emph{Physically Adversarial Infrared Patches with Learnable Shapes and Locations}.
\newblock DOI:10.1109/cvpr52729.2023.01187, 2023.
\newblock \url{https://doi.org/10.1109/cvpr52729.2023.01187}

\bibitem{doi:10.1109/access.2024.3418614}
\newblock \emph{Enhancing Security in Real-Time Video Surveillance: A Deep Learning-Based Remedial Approach for Adversarial Attack Mitigation}.
\newblock DOI:10.1109/access.2024.3418614, 2024.
\newblock \url{https://doi.org/10.1109/access.2024.3418614}

\bibitem{doi:10.1038/s41598-025-15753-8}
\newblock \emph{A comprehensive survey of deep face verification systems adversarial attacks and defense strategies}.
\newblock DOI:10.1038/s41598-025-15753-8, 2025.
\newblock \url{https://doi.org/10.1038/s41598-025-15753-8}

\bibitem{arxiv:1711.00117}
\newblock \emph{Countering Adversarial Images using Input Transformations}.
\newblock arXiv:1711.00117, 2017.
\newblock \url{https://arxiv.org/abs/1711.00117}

\bibitem{arxiv:1911.05268}
\newblock \emph{Adversarial Examples in Modern Machine Learning: A Review}.
\newblock arXiv:1911.05268, 2019.
\newblock \url{https://arxiv.org/abs/1911.05268}

\bibitem{doi:10.1109/sp46215.2023.10179316}
\newblock \emph{On The Empirical Effectiveness of Unrealistic Adversarial Hardening Against Realistic Adversarial Attacks}.
\newblock DOI:10.1109/sp46215.2023.10179316, 2023.
\newblock \url{https://doi.org/10.1109/sp46215.2023.10179316}

\bibitem{arxiv:2408.01541}
\newblock \emph{Guardians of Image Quality: Benchmarking Defenses Against Adversarial Attacks on Image Quality Metrics}.
\newblock arXiv:2408.01541, 2024.
\newblock \url{https://arxiv.org/abs/2408.01541}

\bibitem{doi:10.1109/tpami.2025.3610085}
\newblock \emph{Revisiting Transferable Adversarial Images: Systemization, Evaluation, and New Insights}.
\newblock DOI:10.1109/tpami.2025.3610085, 2025.
\newblock \url{https://doi.org/10.1109/tpami.2025.3610085}

\bibitem{arxiv:2012.12235}
\newblock \emph{Unadversarial Examples: Designing Objects for Robust Vision}.
\newblock arXiv:2012.12235, 2020.
\newblock \url{https://arxiv.org/abs/2012.12235}

\bibitem{doi:10.1609/aaai.v35i9.16927}
\newblock \emph{Attribute-Guided Adversarial Training for Robustness to Natural Perturbations}.
\newblock DOI:10.1609/aaai.v35i9.16927, 2021.
\newblock \url{https://doi.org/10.1609/aaai.v35i9.16927}

\bibitem{doi:10.3389/fnbot.2023.1129720}
\newblock \emph{DualFlow: Generating imperceptible adversarial examples by flow field and normalize flow-based model}.
\newblock DOI:10.3389/fnbot.2023.1129720, 2023.
\newblock \url{https://doi.org/10.3389/fnbot.2023.1129720}

\bibitem{doi:10.3390/cancers15174228}
\newblock \emph{Adversarial Attacks on Medical Image Classification}.
\newblock DOI:10.3390/cancers15174228, 2023.
\newblock \url{https://doi.org/10.3390/cancers15174228}

\bibitem{doi:10.1109/tgrs.2022.3225306}
\newblock \emph{Benchmarking Adversarial Patch Against Aerial Detection}.
\newblock DOI:10.1109/tgrs.2022.3225306, 2022.
\newblock \url{https://doi.org/10.1109/tgrs.2022.3225306}

\bibitem{doi:10.1016/j.jisa.2020.102694}
\newblock \emph{Robust and Natural Physical Adversarial Examples for Object Detectors}.
\newblock DOI:10.1016/j.jisa.2020.102694, 2021.
\newblock \url{https://doi.org/10.1016/j.jisa.2020.102694}

\bibitem{arxiv:2007.16118}
\newblock \emph{Physical Adversarial Attack on Vehicle Detector in the Carla Simulator}.
\newblock arXiv:2007.16118, 2020.
\newblock \url{https://arxiv.org/abs/2007.16118}

\bibitem{doi:10.1109/cvpr52688.2022.01487}
\newblock \emph{DTA: Physical Camouflage Attacks using Differentiable Transformation Network}.
\newblock DOI:10.1109/cvpr52688.2022.01487, 2022.
\newblock \url{https://doi.org/10.1109/cvpr52688.2022.01487}

\bibitem{doi:10.1609/aaai.v34i01.5443}
\newblock \emph{Robust Adversarial Objects against Deep Learning Models}.
\newblock DOI:10.1609/aaai.v34i01.5443, 2020.
\newblock \url{https://doi.org/10.1609/aaai.v34i01.5443}

\bibitem{doi:10.1109/tpami.2022.3176760}
\newblock \emph{Adversarial Sticker: A Stealthy Attack Method in the Physical World}.
\newblock DOI:10.1109/tpami.2022.3176760, 2022.
\newblock \url{https://doi.org/10.1109/tpami.2022.3176760}

\bibitem{doi:10.1109/vr51125.2022.00073}
\newblock \emph{SPAA: Stealthy Projector-based Adversarial Attacks on Deep Image Classifiers}.
\newblock DOI:10.1109/vr51125.2022.00073, 2022.
\newblock \url{https://doi.org/10.1109/vr51125.2022.00073}

\bibitem{doi:10.1016/j.neucom.2022.05.031}
\newblock \emph{Enhancing real-world adversarial patches through 3D modeling of complex target scenes}.
\newblock DOI:10.1016/j.neucom.2022.05.031, 2022.
\newblock \url{https://doi.org/10.1016/j.neucom.2022.05.031}

\bibitem{arxiv:2203.09123}
\newblock \emph{Improving the Transferability of Targeted Adversarial Examples through Object-Based Diverse Input}.
\newblock arXiv:2203.09123, 2022.
\newblock \url{https://arxiv.org/abs/2203.09123}

\bibitem{doi:10.1109/access.2024.3415356}
\newblock \emph{Adversarial Examples for Image Cropping: Gradient-Based and Bayesian-Optimized Approaches for Effective Adversarial Attack}.
\newblock DOI:10.1109/access.2024.3415356, 2024.
\newblock \url{https://doi.org/10.1109/access.2024.3415356}

\bibitem{doi:10.1609/aaai.v38i19.30150}
\newblock \emph{Visual Adversarial Examples Jailbreak Aligned Large Language Models}.
\newblock DOI:10.1609/aaai.v38i19.30150, 2024.
\newblock \url{https://doi.org/10.1609/aaai.v38i19.30150}

\bibitem{doi:10.1109/access.2022.3160283}
\newblock \emph{ARGAN: Adversarially Robust Generative Adversarial Networks for Deep Neural Networks against Adversarial Examples}.
\newblock DOI:10.1109/access.2022.3160283, 2022.
\newblock \url{https://doi.org/10.1109/access.2022.3160283}

\bibitem{doi:10.1145/3459637.3482029}
\newblock \emph{Adversarial Robustness of Deep Learning: Theory, Algorithms, and Applications}.
\newblock DOI:10.1145/3459637.3482029, 2021.
\newblock \url{https://doi.org/10.1145/3459637.3482029}

\bibitem{doi:10.1109/wacv51458.2022.00288}
\newblock \emph{Evaluating the Robustness of Semantic Segmentation for Autonomous Driving against Real-World Adversarial Patch Attacks}.
\newblock DOI:10.1109/wacv51458.2022.00288, 2022.
\newblock \url{https://doi.org/10.1109/wacv51458.2022.00288}

\bibitem{doi:10.1109/iccv.2019.00051}
\newblock \emph{Towards Adversarially Robust Object Detection}.
\newblock DOI:10.1109/iccv.2019.00051, 2019.
\newblock \url{https://doi.org/10.1109/iccv.2019.00051}

\bibitem{doi:10.1016/j.patcog.2022.109009}
\newblock \emph{Robust Physical-World Attacks on Face Recognition}.
\newblock DOI:10.1016/j.patcog.2022.109009, 2022.
\newblock \url{https://doi.org/10.1016/j.patcog.2022.109009}

\bibitem{doi:10.1109/tifs.2023.3278458}
\newblock \emph{The Art of Defense: Letting Networks Fool the Attacker}.
\newblock DOI:10.1109/tifs.2023.3278458, 2023.
\newblock \url{https://doi.org/10.1109/tifs.2023.3278458}

\bibitem{arxiv:1712.02779}
\newblock \emph{Exploring the Landscape of Spatial Robustness}.
\newblock arXiv:1712.02779, 2017.
\newblock \url{https://arxiv.org/abs/1712.02779}

\bibitem{doi:10.1109/icip.2019.8803770}
\newblock \emph{Extending Adversarial Attacks and Defenses to Deep 3D Point Cloud Classifiers}.
\newblock DOI:10.1109/icip.2019.8803770, 2019.
\newblock \url{https://doi.org/10.1109/icip.2019.8803770}

\bibitem{doi:10.1145/3627673.3679750}
\newblock \emph{Towards Robust Vision Transformer via Masked Adaptive Ensemble}.
\newblock DOI:10.1145/3627673.3679750, 2024.
\newblock \url{https://doi.org/10.1145/3627673.3679750}

\bibitem{doi:10.1109/access.2024.3479949}
\newblock \emph{Visual Prompt Engineering for Enhancing Facial Recognition Systems Robustness Against Evasion Attacks}.
\newblock DOI:10.1109/access.2024.3479949, 2024.
\newblock \url{https://doi.org/10.1109/access.2024.3479949}

\bibitem{doi:10.1145/3449726.3463137}
\newblock \emph{Black-box adversarial attacks using evolution strategies}.
\newblock DOI:10.1145/3449726.3463137, 2021.
\newblock \url{https://doi.org/10.1145/3449726.3463137}

\bibitem{doi:10.1109/satml54575.2023.00030}
\newblock \emph{Towards Transferable Unrestricted Adversarial Examples with Minimum Changes}.
\newblock DOI:10.1109/satml54575.2023.00030, 2023.
\newblock \url{https://doi.org/10.1109/satml54575.2023.00030}

\bibitem{doi:10.1109/access.2024.3446834}
\newblock \emph{Uncovering Distortion Differences: A Study of Adversarial Attacks and Machine Discriminability}.
\newblock DOI:10.1109/access.2024.3446834, 2024.
\newblock \url{https://doi.org/10.1109/access.2024.3446834}

\bibitem{arxiv:2301.10766}
\newblock \emph{On the Adversarial Robustness of Camera-based 3D Object Detection}.
\newblock arXiv:2301.10766, 2023.
\newblock \url{https://arxiv.org/abs/2301.10766}

\bibitem{doi:10.1109/icpr48806.2021.9412236}
\newblock \emph{AdvHat: Real-World Adversarial Attack on ArcFace Face ID System}.
\newblock DOI:10.1109/icpr48806.2021.9412236, 2021.
\newblock \url{https://doi.org/10.1109/icpr48806.2021.9412236}

\bibitem{doi:10.1109/jstars.2024.3422377}
\newblock \emph{Pattern Corruption-Assisted Physical Attacks Against Object Detection in UAV Remote Sensing}.
\newblock DOI:10.1109/jstars.2024.3422377, 2024.
\newblock \url{https://doi.org/10.1109/jstars.2024.3422377}

\bibitem{doi:10.1109/vrw62533.2024.00284}
\newblock \emph{Diffusion Attack: Leveraging Stable Diffusion for Naturalistic Image Attacking}.
\newblock DOI:10.1109/vrw62533.2024.00284, 2024.
\newblock \url{https://doi.org/10.1109/vrw62533.2024.00284}

\bibitem{doi:10.59247/csol.v2i1.73}
\newblock \emph{A Review on Attacks against Artificial Intelligence (AI) and Their Defence Image Recognition and Generation Machine Learning, Artificial Intelligence}.
\newblock DOI:10.59247/csol.v2i1.73, 2024.
\newblock \url{https://doi.org/10.59247/csol.v2i1.73}

\end{thebibliography}

\end{document}
